\chapter{Anesthesiology}

\medterm{Acetaminophen (Paracetamol)}-- A widely used non-opioid, non-NSAID analgesic and antipyretic agent that forms the foundational first step of multimodal perioperative analgesia. Its primary mechanism involves central inhibition of cyclooxygenase (COX-1 and COX-2) enzymes with selectivity for peroxidase sites, facilitation of descending serotonergic pain-inhibitory pathways, and active cannabinoid/TRPA1 receptor metabolite interactions (AM404). Unlike traditional NSAIDs, it possesses minimal peripheral anti-inflammatory action, does not impair platelet aggregation, does not damage gastric mucosa, and does not alter renal blood flow. Administered intravenously (1000 mg IV over 15 minutes every 6 hours; pediatric dose 15 mg/kg; maximum 4000 mg/day in adults, or 2000--3000 mg/day in chronic alcohol misuse, malnutrition, or hepatic impairment). Provides a significant opioid-sparing effect (reducing postoperative opioid consumption by 15 to 30 percent). Metabolism is predominantly hepatic via glucuronidation and sulfation; a minor fraction is metabolized by CYP2E1 into the reactive, hepatotoxic intermediate N-acetyl-p-benzoquinone imine (NAPQI), which is normally detoxified by glutathione. Overdose depletes glutathione, causing acute hepatic centrilobular necrosis, treated specifically with intravenous N-acetylcysteine (NAC).

\medterm{Airway Management}-- The systematic evaluation, preparation, and execution of techniques and devices to maintain patent ventilation, oxygenation, and airway protection. Preoperative airway assessment utilizes clinical predictors of difficult intubation, including the modified Mallampati classification: Class I (soft palate, fauces, uvula, and tonsillar pillars visible), Class II (soft palate, fauces, and uvula visible), Class III (soft palate and base of uvula visible), and Class IV (only hard palate visible). Other critical predictors include inter-incisor mouth opening (<3 cm), thyromental distance (<6.5 cm), sternomental distance, neck circumference (>43 cm), limited cervical spine mobility (<35 degrees extension), and the upper lip bite test. A designated difficult airway cart—equipped with video laryngoscopes, supraglottic airway devices, flexible bronchoscopes, gum-elastic bougies, and scalpel-cricothyrotomy emergency kits—must be immediately accessible whenever airway difficulty is anticipated.

\medterm{Anaphylaxis During Anesthesia}-- A life-threatening systemic hypersensitivity reaction that may cause sudden hypotension, bronchospasm, cardiovascular collapse, and sometimes urticaria or flushing. Common perioperative triggers include neuromuscular-blocking agents, antibiotics, chlorhexidine, and latex. A sudden fall in end-tidal carbon dioxide may accompany severe reduced pulmonary blood flow. Management requires stopping the suspected trigger, calling for help, administering high-concentration oxygen, supporting the airway and circulation, giving titrated intravenous or intramuscular epinephrine according to the clinical setting, and providing intravenous fluid resuscitation; antihistamines and corticosteroids are adjuncts, not substitutes for epinephrine.

\medterm{Anesthesia}-- The use of medications to prevent sensation, particularly pain, during medical procedures or surgery. Types: general anesthesia (reversible loss of consciousness and sensation, involving hypnotics, analgesics, and muscle relaxants), regional anesthesia (blocking sensation in a specific body region, e.g., epidural, spinal, brachial plexus block), local anesthesia (numbing a small area, e.g., infiltration, topical), and sedation (reduced consciousness while maintaining airway reflexes and spontaneous ventilation). General anesthesia involves three components: hypnosis (loss of consciousness, propofol, volatile anesthetics), analgesia (pain relief, opioids, ketamine, lidocaine), and muscle relaxation (neuromuscular blocking agents). Monitoring: ECG, BP, SpO$_2$, end-tidal CO$_2$, anesthetic gas concentration, and temperature.

\medterm{Anesthesia for Coagulopathy}-- Coagulopathy increases bleeding risk during surgery. Assessment includes history, physical examination, and laboratory tests (INR, PT, aPTT, platelet count, fibrinogen). INR greater than 1.5 or aPTT greater than 1.5 times control increases bleeding risk. Platelet count less than 50,000 for major surgery or less than 100,000 for neurosurgery requires platelet transfusion. Fibrinogen less than 150 mg/dL requires cryoprecipitate or fibrinogen concentrate. Fresh frozen plasma (15 to 20 mL/kg) corrects coagulopathy effectively. Vitamin K (10 mg IV slow push, 30 minutes onset) is used for warfarin reversal. Prothrombin complex concentrate (PCC) is preferred over FFP for urgent reversal (INR correction in 15 minutes). Antifibrinolytics (tranexamic acid 1 g IV) reduce bleeding in trauma and cardiac surgery.

\medterm{Anesthesia for Epilepsy}-- Patients with epilepsy generally require continuation of antiseizure medication perioperatively because missed doses increase seizure risk. Propofol and benzodiazepines are anticonvulsant; high concentrations of some volatile agents may produce epileptiform activity. Hyperventilation, hypoglycemia, electrolyte abnormalities, sleep deprivation, and inadequate analgesia may precipitate seizures. Meperidine and tramadol may lower the seizure threshold. Convulsive status epilepticus requires immediate airway, breathing, and circulation support, rapid benzodiazepine treatment, and escalation to a longer-acting antiseizure medication if seizures persist.

\medterm{Anesthesia for Hemophilia}-- Hemophilia A results from factor VIII deficiency and hemophilia B from factor IX deficiency. Perioperative care is coordinated with hematology and uses factor replacement or, in selected mild hemophilia A cases, desmopressin. Avoid unnecessary traumatic or intramuscular procedures, and perform neuraxial blockade only when coagulation status and factor activity meet accepted specialist thresholds. Monitor factor activity and clinical bleeding before and after surgery.

\medterm{Anesthesia for Neurosurgery}-- Neuroanesthesia aims to maintain cerebral perfusion pressure (CPP = mean arterial pressure minus intracranial pressure), optimize brain relaxation, and minimize secondary brain injury. Induction agents (propofol, etomidate) reduce cerebral metabolic rate and intracranial pressure. Volatile anesthetics increase cerebral blood flow at concentrations greater than 1 MAC. Hyperventilation (PaCO2 30 to 35 mmHg) causes cerebral vasoconstriction and reduces intracranial pressure acutely. Mannitol (0.5 to 1 g/kg IV) reduces intracranial pressure through osmotic diuresis. Maintain CPP at 60 to 70 mmHg. Avoid hypotension and hypoxia. Emergence should be smooth to prevent coughing and hypertension. Postoperative neurological assessment is essential.

\medterm{Anesthesia for Pheochromocytoma}-- Anesthesia for patients with pheochromocytoma (a rare, catecholamine-secreting neuroendocrine tumor of chromaffin cells) represents a high-risk scenario characterized by severe hemodynamic instability. Preoperative medical preparation is vital and follows Roizen criteria: initiation of non-competitive alpha-adrenergic blockade (phenoxybenzamine 10--14 days preoperatively, titrated until orthostatic hypotension is tolerated) or selective alpha-1 blockers (doxazosin), followed strictly by beta-adrenergic blockade (propranolol, labetalol) only AFTER adequate alpha-blockade has been established to prevent catastrophic unopposed alpha-mediated vasoconstriction and hypertensive crisis. Intraoperative hemodynamic monitoring requires an arterial line placed prior to induction and large-bore central venous access. Induction, laryngoscopy, pneumoperitoneum, and surgical manipulation of the tumor provoke massive surges of epinephrine and norepinephrine, managed with rapid-acting intravenous vasodilators and beta-blockers (sodium nitroprusside, nicardipine, phentolamine, esmolol, magnesium sulfate). Crucially, immediately following tumor ligation and resection, circulating catecholamine levels plummet precipitously, causing severe refractory hypotension due to down-regulated vascular adrenergic receptors and residual preoperative alpha-blockade. Continued postoperative alpha and beta blockade is strictly contraindicated; management requires aggressive volume expansion, vasopressor support (norepinephrine, vasopressin), and close monitoring for rebound hypoglycemia resulting from sudden uninhibited insulin secretion.

\medterm{Anesthesia for Porphyria}-- Acute intermittent porphyria is a metabolic disorder of heme synthesis causing abdominal pain, neuropathy, psychiatric symptoms, and dark urine. Certain anesthetics are porphyrinogenic (triggering attacks): barbiturates, sulfonamides, carbamazepine, and possibly volatile anesthetics (though sevoflurane is considered safe). Safe agents include propofol, opioids, benzodiazepines (low dose), succinylcholine, and non-depolarizing neuromuscular blockers. Porphyria-safe drug databases should be consulted and porphyria-safe drug lists verified before administering any anesthetic agent.

\medterm{Anesthesia Gas Analyzers}-- Anesthesia gas analyzers measure delivered and exhaled concentrations of oxygen, carbon dioxide, and volatile anesthetic agents. The oxygen analyzer measures inspired oxygen concentration; pulse oximetry measures peripheral hemoglobin oxygen saturation. Capnography measures exhaled carbon dioxide and helps confirm tracheal intubation, assess ventilation, and detect disconnection or obstruction. Temperature and neuromuscular monitoring provide additional physiologic safety information. Arterial pressure, central venous pressure, and cardiac output monitoring are selected according to the patient's condition and procedure.

\medterm{Anesthesia Informatics}-- Anesthesia informatics integrates electronic health records, decision support systems, and data analytics into clinical practice. Automated documentation reduces charting time. Clinical decision support alerts for drug interactions and dosing errors. Big data analytics identify patterns in anesthesia outcomes.

\medterm{Anesthesia Machine}-- The anesthesia machine delivers precise mixtures of oxygen, air, and volatile anesthetics to the patient. Components include gas cylinders (E-cylinder oxygen contains 660 liters at 2000 psi), pipeline supplies (hospital oxygen at 50 psi), pressure regulators, fail-safe valve (oxygen pressure below 25 psi shuts off nitrous oxide), vaporizers (variable-bypass sevoflurane, desflurane uses heated pressurized system), scavenging system (removes waste anesthetic gases), and breathing circuit (Circle system) comprises fresh gas inlet, inspiratory and expiratory valves, reservoir bag, CO2 absorber, adjustable pressure limiting valve, and ventilator. The circle system reduces gas consumption and heat loss. Safety features include oxygen failure alarm, hypoxia prevention system, and vapouriser interlock (prevents simultaneous use of two volatile agents).

\medterm{Anesthesia Monitoring}-- Anesthesia monitoring encompasses continuous assessment of patient physiology during surgery. Standard ASA monitoring includes electrocardiography (detecting arrhythmias and ischemia), pulse oximetry (continuous SpO2 measurement), non-invasive blood pressure (oscillometric every 3 to 5 minutes), end-tidal CO2 (confirming tube position and ventilation adequacy), temperature monitoring, and inspired oxygen concentration. Invasive arterial pressure monitoring provides beat-to-beat blood pressure and allows precise vasopressor management. Central venous pressure and cardiac output monitoring are used in high-risk patients. Neuromuscular monitoring (train-of-four) is used for muscle relaxant management. Temperature monitoring is essential for detecting malignant hyperthermia and hypothermia. The goal is to maintain physiological parameters within acceptable ranges and detect abnormalities early.

\medterm{Anesthesia Pharmacogenomics}-- Anesthesia pharmacogenomics studies genetic variations affecting drug response. Pseudocholinesterase deficiency causes prolonged succinylcholine paralysis. CYP2D6 polymorphisms affect codeine metabolism to morphine. Malignant hyperthermia susceptibility is linked to RYR1 gene mutations. Pharmacogenomic testing guides personalized anesthesia.

\medterm{Arterial Catheterization}-- The placement of an indwelling catheter into an accessible peripheral artery to enable continuous, real-time beat-to-beat arterial blood pressure monitoring and frequent blood sampling for arterial blood gas (ABG), electrolyte, and lactate analysis. The radial artery at the wrist is the most frequent insertion site due to its superficial course and dual collateral blood supply to the hand via the ulnar artery (historically evaluated by the modified Allen test, though ultrasound evaluation is now standard); alternative sites include the femoral, axillary, brachial, and dorsalis pedis arteries. The catheter is connected via non-compliant, fluid-filled tubing to an electromechanical pressure transducer leveled at the phlebostatic axis (fourth intercostal space, mid-axillary line, representing the right atrium) and zeroed to atmospheric pressure. Dynamic arterial waveform analysis yields physiological parameters including pulse pressure variation (PPV) and stroke volume variation (SVV) for functional fluid responsiveness in mechanically ventilated patients. Complications include catheter-related thrombosis, arterial spasm, distal digital ischemia, hematoma formation, pseudoaneurysm, and catheter-related bloodstream infection.

\medterm{ASA Physical Status Classification System}-- A system developed by the American Society of Anesthesiologists to assess and communicate a patient's pre-operative physical status and co-morbidities. It ranges from ASA I (a normal healthy patient) to ASA VI (a declared brain-dead patient whose organs are being removed for donor purposes). An 'E' is appended to the classification (e.g., ASA IIIE) if the surgery is an emergency. It is a key tool for risk stratification and clinical documentation.

\medterm{Aspiration Pneumonitis}-- A chemical lung injury (historically known as Mendelson syndrome) resulting from the inhalation of sterile, acidic gastric contents into the pulmonary tree. Severe chemical pneumonitis occurs classically when aspirated gastric volume exceeds 0.4 mL/kg (approx. 25 mL in adults) with a pH below 2.5. The acidic insult causes immediate desquamation of bronchial epithelium, destruction of type I and II alveolar pneumocytes, interstitial edema, and severe chemical inflammatory response within 1--2 hours, presenting as bronchospasm, tachypnea, cyanosis, acute hypoxemia, and patchy infiltrates on chest radiography. Differentiated from aspiration pneumonia, which is an infectious process resulting from aspiration of colonized oropharyngeal secretions. Immediate management involves rapid suctioning of the oropharynx, endotracheal intubation, positive end-expiratory pressure (PEEP), bronchodilators for bronchospasm, and lung-protective mechanical ventilation; prophylactic antibiotics and corticosteroids are not recommended.

\medterm{Autologous Blood Donation}-- Autologous blood donation involves preoperative collection and storage of the patient's own blood for intraoperative or postoperative transfusion. It eliminates risks of alloimmune reactions, transfusion-transmitted infections, and immune modulation. Indications include rare blood types, multiple alloantibodies, elective orthopedic surgery, and religious objections to allogeneic transfusion. Preoperative donation requires hemoglobin greater than 11 g/dL and at least 3 weeks before surgery. Intraoperative cell salvage and acute normovolaemic haemodilution reduce allogeneic transfusion requirements. Complications include bacterial contamination, clerical error, and fluid overload.

\medterm{Awake Fiberoptic Intubation}-- The gold standard for anticipated difficult airways, performed with the patient conscious and cooperatively maintaining spontaneous ventilation. Topical anesthesia is achieved through lidocaine nebulization, nerve blocks (glossopharyngeal, superior laryngeal, transtracheal), and spray-as-you-go technique. Sedation with dexmedetomidine or remifentanil provides anxiolysis while maintaining respiratory drive. The fiberoptic bronchoscope is passed through the vocal cords, and the endotracheal tube is passed over the bronchoscope into the trachea. Awake fiberoptic intubation is used in patients with known difficult airway, limited neck mobility (cervical spine injury, rheumatoid arthritis), head and neck tumours, and morbid obesity. Contraindications include patient refusal, severe hypoxia, and bleeding diathesis.

\medterm{Awareness Under Anesthesia}-- Unintended recall of events during general anesthesia, most commonly occurring during procedures where neuromuscular blockade prevents movement. Incidence: 0.1-0.2\% of all general anesthetics, higher in cardiac, trauma, and cesarean section. Presents with postoperative recall of conversations, surgical sounds, or pain (worst outcomes). Risk factors: inadequate anesthesia (planned light anesthesia, hemodynamic instability), equipment failure, and certain drugs (ketamine, etomidate, nitrous oxide). Prevention: adequate dosing, volatile anesthetic end-tidal monitoring (MAC >0.7), BIS monitoring (target 40-60), and pre-oxygenation with a protocol. Management: documented debriefing, apology, psychological support, and referral to post-traumatic stress disorder services.

\medterm{Blood Products}-- Therapeutic biological components derived from whole blood fractionation, administered during surgery to restore circulating volume, oxygen-carrying capacity, and hemostatic function. Components include: (1) Packed Red Blood Cells (PRBCs, hematocrit ~55--65\%), indicated to restore oxygen transport in anemia or acute blood loss; 1 unit raises adult hemoglobin by ~1 g/dL. (2) Fresh Frozen Plasma (FFP), containing all plasma coagulation factors and physiological anticoagulants, indicated for multiple factor deficiencies, severe liver disease, warfarin reversal when PCC is unavailable, and massive transfusion; standard dose 10--15 mL/kg. (3) Platelets (apheresis unit or pooled concentrate), indicated for thrombocytopenia ($<50,000/\mu$L for major surgery, $<100,000/\mu$L for neurosurgery) or platelet dysfunction; 1 apheresis unit increases count by 30,000--50,000/$\mu$L. (4) Cryoprecipitate, prepared by thawing FFP, rich in fibrinogen ($>150$ mg/unit), factor VIII, von Willebrand factor, and factor XIII, used for hypofibrinogenemia ($<150--200$ mg/dL) and massive hemorrhage. (5) Whole Blood, increasingly used in military and trauma resuscitation for massive hemorrhage to provide physiological reconstitution in a single product. Transfusion reactions include acute hemolytic reactions (ABO incompatibility), febrile non-hemolytic reactions, allergic/anaphylactic reactions, transfusion-related acute lung injury (TRALI), transfusion-associated circulatory overload (TACO), and transmission of viral/bacterial pathogens.

\medterm{Brain Monitoring}-- The intraoperative assessment of cerebral electrophysiology, depth of hypnosis, and neurological integrity during general anesthesia. Processed electroencephalography (pEEG), such as the Bispectral Index (BIS), evaluates frontal raw EEG signals via Fourier transform and spectral analysis, generating a dimensionless index from 0 (isoelectric/cerebral silence) to 100 (fully awake). A target BIS range of 40 to 60 is recommended for general anesthesia to minimize the risk of unintended intraoperative awareness while avoiding excessively deep anesthesia linked to postoperative delirium and mortality. Similar devices include Patient State Index (SedLine) and Spectral Entropy (State and Response Entropy). Limitations: pEEG values can be paradoxical or unreliable in the presence of ketamine and nitrous oxide (which produce high-frequency frontal EEG activity without awakening), hypothermia, severe encephalopathy, or electromyographic artifact. In vascular and neurosurgery, continuous raw EEG and evoked potentials provide sensitive real-time detection of cerebral ischemia and cortical hypoperfusion.

\medterm{Bronchospasm}-- Involuntary constriction of bronchial smooth muscle causing airway obstruction, wheezing, and increased airway pressures. Triggers include airway manipulation, volatile anesthetics (sevoflurane is bronchodilating; desflurane can cause bronchospasm), histamine-releasing drugs (morphine, atracurium), and reactive airway disease. Treatment includes deepening anesthesia with sevoflurane, beta-2 agonists (albuterol nebulized), epinephrine (10 to 100 mcg IV), and corticosteroids (hydrocortisone 100 mg IV). Volatile anesthetics may need to be reduced or switched to intravenous agents (propofol) to permit effective bronchodilator therapy, and severe cases may require intravenous adrenaline infusion and magnesium sulphate (2 g IV).

\medterm{Bupivacaine}-- A potent, long-acting amino-amide local anesthetic widely utilized for infiltration, peripheral nerve blocks, epidural analgesia/anesthesia, and spinal anesthesia. It exists as a racemic mixture of S(-) and R(+) enantiomers. It blocks voltage-gated sodium channels in nerve membranes, preventing action potential propagation. Because of its high lipid solubility and strong protein binding (approximately 95 percent, primarily to alpha-1 acid glycoprotein), it provides a prolonged duration of anesthesia (3 to 6 hours for nerve blocks). It exhibits marked sensory-motor differentiation at lower concentrations (0.0625--0.125\%), making it the historical gold standard for obstetric labor analgesia; higher concentrations (0.5--0.75\%) provide dense surgical motor blockade. Hyperbaric bupivacaine (0.75\% in 8.25\% dextrose) is the most common agent for subarachnoid spinal anesthesia. Maximum safe dose is 2.0 to 2.5 mg/kg without epinephrine (maximum 175 mg), and up to 3.0 mg/kg with epinephrine (maximum 225 mg). Major hazard: Bupivacaine is the most cardiotoxic local anesthetic in routine use because the R(+) enantiomer binds avidly to and dissociates extremely slowly from cardiac sodium channels, predisposing to refractory re-entrant ventricular arrhythmias, electromechanical dissociation, and cardiac arrest in systemic toxicity (LAST), requiring resuscitation with 20\% lipid emulsion (Intralipid).

\medterm{Cannot Intubate Cannot Oxygenate}-- The most critical airway emergency — inability to intubate the trachea and inability to oxygenate the patient by face mask or supraglottic airway despite optimal efforts. CICO is a life-threatening situation requiring immediate surgical airway access. Signs: SpO2 decreasing despite 100\% O2, absent end-tidal CO2, no chest rise. Management (CICO algorithm): declare CICO, call for surgical airway, perform scalpel cricothyroidotomy (vertical skin incision, horizontal membrane incision, tracheal hook, insert endotracheal tube or bougie). Alternative: needle cricothyroidotomy with jet ventilation (temporizing). Prevention: early recognition of difficult airway, awake techniques, and capnography confirmation of every intubation.

\medterm{Capnography}-- The continuous, noninvasive real-time measurement and graphical display of carbon dioxide concentration in the breathing circuit throughout the respiratory cycle, providing immediate information on ventilation, pulmonary perfusion, and cellular metabolism. Utilizes infrared spectroscopy to measure partial pressure of CO2 (PCO2). The normal capnogram exhibits four distinct phases: Phase I (inspiratory baseline, dead space gas with zero CO2); Phase II (rapid expiratory upstroke, representing mixed anatomical dead space and alveolar gas); Phase III (alveolar plateau, representing exhalation of pure alveolar gas, terminating at the highest point: the end-tidal CO2 [$EtCO_2$]); and Phase IV (rapid downstroke to zero as inspiration begins). Normal $EtCO_2$ is 35 to 45 mmHg, which typically reflects arterial PCO2 ($PaCO_2$) with an expected arterial-to-end-tidal gradient of 2 to 5 mmHg in healthy lungs. Clinical applications: (1) Immediate confirmation of endotracheal tube placement (presence of $\ge 6$ continuous, consistent capnographic waveforms is the gold standard for ruling out esophageal intubation); (2) Assessment of effective chest compressions and early indicator of return of spontaneous circulation (ROSC) during CPR; (3) Early detection of malignant hyperthermia (exponential rise in $EtCO_2$); (4) Detection of pulmonary embolism, hypovolemia, or cardiac arrest (sudden profound drop in $EtCO_2$ due to increased alveolar dead space); and (5) Identification of obstructive airway disease (prolonged 'shark-fin' upstroke in Phase III caused by uneven expiratory flow in asthma or COPD).

\medterm{Carbon Dioxide Absorber}-- A safety device used in anesthesia breathing circuits that filters out carbon dioxide from a patient's exhaled breath, allowing the remaining anesthetic gases to be safely recycled and rebreathed.

\medterm{Central Venous Catheter}-- An indwelling vascular access catheter whose tip resides in a great intrathoracic vein, typically the lower third of the superior vena cava (SVC) or the cavoatrial junction. Standard insertion sites include the right internal jugular vein (preferred due to straight trajectory into the SVC and low complication rate under real-time ultrasound guidance), the subclavian vein (associated with the lowest infection rate and patient comfort, but higher pneumothorax risk), and the femoral vein (useful in emergencies, but higher deep vein thrombosis and infection risk). Indications include: administration of caustic, sclerosing, or vasoactive infusions (e.g., norepinephrine, vasopressin, concentrated potassium, hypertonic solutions), total parenteral nutrition (TPN), central venous pressure (CVP) monitoring, transvenous cardiac pacing, hemodialysis/plasmapheresis, and venous access when peripheral veins are exhausted. Placement is performed standardly using the Seldinger wire-guided technique with ultrasound visualization. Complications include accidental arterial puncture, hematoma, pneumothorax/hemothorax, air embolism, cardiac arrhythmias (catheter tip irritating right atrium), catheter-related bloodstream infection (CRBSI), and central venous thrombosis.

\medterm{Chronic Pain}-- Pain persisting longer than three months or beyond normal tissue healing time, often lacking an ongoing identifiable cause. It involves complex interactions between biological, psychological, and social factors that can amplify pain signaling in the nervous system. Management is multimodal, combining medications, physical therapy, cognitive behavioral therapy, nerve blocks, and complementary approaches.

\medterm{Chronic Pain Management}-- Chronic pain management extends beyond the perioperative period. Opioid therapy for chronic pain requires careful patient selection, informed consent, and regular reassessment. Multimodal approach includes NSAIDs, acetaminophen, gabapentinoids, antidepressants (SNRIs, TCAs), topical agents (lidocaine patches, capsaicin), and interventional procedures (nerve blocks, radiofrequency ablation, spinal cord stimulation). Opioid rotation (switching to a different opioid) is used for tolerance or side effects; methadone is used for both pain and opioid use disorder due to its long half-life (24-36 hours) and NMDA receptor antagonism, which may reduce tolerance. Buprenorphine is a partial mu-opioid agonist with ceiling effect on respiratory depression, making it safer in overdose but with a theoretical risk of precipitating withdrawal. Adjuvant therapies include physical therapy, cognitive behavioural therapy, and neuromodulation (spinal cord stimulation, peripheral nerve stimulation). Long-term opioid use carries risks of hypogonadism, immunosuppression, and opioid-induced hyperalgesia.

\medterm{Cisatracurium}-- An intermediate-acting non-depolarizing neuromuscular blocking agent, representing one of the ten stereoisomers of atracurium (the 1R-cis, 1'R-cis isomer), possessing approximately three to four times the neuromuscular blocking potency of the parent drug. It competitively blocks nicotinic acetylcholine receptors at the skeletal muscle motor endplate. Standard intubation dose is 0.15 to 0.2 mg/kg IV (onset 2 to 3 minutes, clinical duration 45 to 60 minutes); maintenance infusion is 1 to 2 mcg/kg/min. Its defining clinical feature is its organ-independent elimination: it undergoes spontaneous chemical degradation at normal body temperature and blood pH known as Hofmann elimination (accounting for ~80 percent of clearance) and plasma ester hydrolysis. Because it does not rely on renal or hepatic excretion, its duration of action and recovery profile are completely unaffected by acute renal failure or hepatic cirrhosis, making it the neuromuscular blocker of choice for patients with renal failure, liver failure, and critically ill patients undergoing prolonged infusions in the intensive care unit. Unlike atracurium, cisatracurium does not trigger dose-dependent histamine release, causing no hypotension, tachycardia, or bronchospasm even at large bolus doses.

\medterm{Colloids}-- Intravenous fluids containing high-molecular-weight oncotic macromolecules suspended in an electrolyte carrier solution that remain predominantly within the intravascular compartment, exerting oncotic pressure to expand plasma volume. Natural colloids include human albumin (5\% iso-oncotic, providing ~1:1 volume expansion; and 25\% hyperoncotic, drawing fluid from the interstitial space to expand intravascular volume by 4 to 5 times the infused volume), utilized for large-volume paracentesis in cirrhosis, spontaneous bacterial peritonitis, nephrotic syndrome, and major burn resuscitation. Synthetic colloids include hydroxyethyl starch (HES, derived from waxy maize or potato starch), dextrans, and gelatin solutions. Large clinical trials (CHEST, 6S, SAFE) demonstrated that synthetic starches offer no survival benefit over crystalloids and are associated with acute kidney injury requiring renal replacement therapy, coagulopathy (interference with factor VIII and von Willebrand factor), and pruritus, leading regulatory agencies (FDA, EMA) to restrict their clinical use. Dextran (40 and 70 kDa) is associated with severe anaphylactoid reactions, interference with blood cross-matching, and impaired platelet function.

\medterm{Cricoid Pressure}-- A maneuver (Sellick maneuver) designed to reduce the risk of regurgitation and pulmonary aspiration of gastric contents during rapid sequence induction (RSI) of general anesthesia. It involves the application of backward pressure on the cricoid cartilage with the thumb and index finger, pressing the complete cartilaginous ring posteriorly against the bodies of the cervical vertebrae (C6) to occlude the upper esophagus. Standard clinical protocols recommend applying approximately 10 N (approx. 1 kg) of force in the awake patient to avoid discomfort, retching, or coughing, increasing to 30 N (approx. 3 kg) of force upon loss of consciousness. Cricoid pressure should be maintained continuously until the endotracheal tube cuff is properly inflated, correct position is verified by capnography, and the airway is definitively secured. The maneuver should be released immediately if active vomiting occurs (to prevent esophageal rupture) or if it interferes with bag-mask ventilation, supraglottic airway seating, or direct/video laryngoscopic view. Current difficult airway guidelines emphasize individualizing its use and releasing pressure promptly if glottic visualization is impaired.

\medterm{Cricothyrotomy}-- An emergent surgical airway procedure performed to establish immediate subglottic oxygenation and ventilation in a 'Cannot Intubate, Cannot Oxygenate' (CICO) crisis when all non-invasive airway techniques have failed. The anatomical landmark is the cricothyroid membrane, situated between the inferior border of the thyroid cartilage and the superior border of the cricoid cartilage. The Difficult Airway Society (DAS) recommended technique is the 'Scalpel-Bougie-Tube' method: (1) In a fully extended neck, the cricothyroid membrane is palpated; (2) A transverse stab incision is made through the skin and cricothyroid membrane using a scalpel with a No. 10 or No. 20 blade; (3) The blade is rotated 90 degrees with the sharp edge facing caudally to keep the stoma patent; (4) A coudé-tipped gum-elastic bougie is passed through the incision into the trachea to a depth of 10--15 cm; and (5) A cuffed endotracheal tube (typically size 5.0 to 6.0 mm internal diameter) is railroaded over the bougie into the trachea, the cuff is inflated, and ventilation with 100\% oxygen is initiated, confirmed immediately by capnography. Needle cricothyrotomy with high-pressure transtracheal jet ventilation is an alternative temporizing measure but carries high failure rates and barotrauma risks.

\medterm{Crystalloids}-- Aqueous solutions of mineral salts that distribute between intravascular and interstitial compartments. Normal saline (0.9 percent NaCl) has 154 mEq/L each of sodium and chloride, pH 5.4, and is associated with hyperchloremic metabolic acidosis at large volumes. Lactated Ringer's contains sodium 130, potassium 4, calcium 3, chloride 109, and lactate 28 mEq/L, with pH 6.5. Lactate is metabolized to bicarbonate, making it physiologically balanced. Balanced crystalloids are preferred over normal saline for large-volume resuscitation (>2 L) to reduce the risk of hyperchloraemic metabolic acidosis and acute kidney injury, though isotonic saline remains the most widely available crystalloid worldwide and is appropriate for most clinical situations.

\medterm{Delayed Emergence}-- Prolonged time to regain consciousness after discontinuation of anesthetic agents. Failure to regain consciousness within 30-60 minutes of completing surgery. Causes: anesthetic drug overdose, residual volatile anesthetic (closed circuit, long duration), synergistic sedative effects (benzodiazepines, opioids), residual neuromuscular blockade, metabolic (hypoglycemia, hyponatremia, hypothermia, hypoxia, hypercarbia), and CNS injury (stroke, cerebral edema, seizure). Workup: ABCs, check TOF, end-tidal gas, temperature, glucose, electrolytes, and arterial blood gas. CT head and EEG for suspected neurologic injury.

\medterm{Delirium in Anesthesia}-- Delirium is an acute confusional state characterized by fluctuating consciousness, inattention, disorganized thinking, and altered level of awareness. Postoperative delirium occurs in 15 to 50 percent of elderly surgical patients, with higher rates in hip fracture and cardiac surgery. Risk factors include advanced age, preexisting cognitive impairment, sensory deprivation, sleep disruption, pain, opioids, anticholinergic drugs, and metabolic derangements. Delirium increases length of stay, complications, and mortality; prevention focuses on early recognition of at-risk patients, optimisation of pain and sleep, minimising use of anticholinergic medications, providing sensory aids (glasses, hearing aids), early mobilisation, and family involvement in care.

\medterm{Desflurane}-- A highly fluorinated methyl ethyl ether volatile inhalational anesthetic characterized by the lowest blood-gas partition coefficient (0.42) of all currently available potent inhalational agents. Because of its exceptionally low solubility in blood and body tissues, alveolar concentration rises rapidly toward inspired concentration during induction, and washes out with unmatched speed upon discontinuation, resulting in extremely rapid, clear-headed emergence that is virtually independent of anesthetic duration. Its minimum alveolar concentration (MAC) is approximately 6.0\% in pure oxygen in young adults, making it the least potent volatile agent. Its high vapor pressure (669 mmHg at 20°C, boiling point 22.8°C) requires administration via a dedicated, electrically heated and pressurized vaporizer (Tec 6 or equivalent) maintained at 39°C and 2 atmospheres pressure. Major clinical limitations include: (1) Marked airway pungency and irritation, which reliably induces coughing, breath-holding, copious secretions, and laryngospasm on inhalational induction (precluding its use for mask induction); and (2) Transient sympathetic hyperactivity: rapid increases in inspired desflurane concentration provoke sudden surges in central sympathetic outflow, producing pronounced hypertension and tachycardia.

\medterm{Dexmedetomidine}-- A highly potent, selective alpha-2 adrenergic receptor agonist with an alpha-2 to alpha-1 selectivity ratio of approximately 1620:1 (roughly eight times more selective than clonidine). It acts on presynaptic alpha-2 receptors in the locus coeruleus of the brainstem to inhibit norepinephrine release, producing a unique state of 'cooperative' or 'conscious' sedation that closely mimics non-REM natural sleep, allowing patients to be easily aroused by verbal or tactile stimuli while preserving spontaneous ventilation and upper airway patency without respiratory depression. It also acts on alpha-2 receptors in the spinal cord dorsal horn to provide moderate analgesia and opioid-sparing effects. Clinical dosing: Initial loading infusion of 0.5 to 1.0 mcg/kg IV over 10 to 15 minutes, followed by maintenance infusion of 0.2 to 0.7 mcg/kg/hr. Commonly used for sedation during awake fiberoptic intubation, procedural sedation, postoperative ICU mechanical ventilation weaning, and prevention of emergence delirium in pediatric patients. Adverse effects include biphasic hemodynamic responses: initial transient hypertension (due to peripheral postjunctional alpha-2B vasoconstriction with rapid bolus) followed by persistent dose-dependent bradycardia and hypotension (due to central sympathetic outflow inhibition).

\medterm{Difficult Airway}-- A clinical situation in which a trained anesthesiologist or emergency physician encounters difficulty with bag-mask ventilation, supraglottic airway placement, laryngoscopy, or endotracheal intubation. Difficult mask ventilation: inability to maintain SpO2 >92\% using 100\% oxygen and positive pressure with a mask. Difficult laryngoscopy: Cormack-Lehane grade III or IV view (only epiglottis or no glottic structures visible). Difficult intubation: >3 attempts or >10 minutes required. Predictors: LEMON (Look, Evaluate, Mallampati, Obstruction, Neck mobility), 3-3-2 rule, upper lip bite test, obesity, and OSA. Management: ASA difficult airway algorithm; call for help, optimize position, use video laryngoscope, bougie, and consider awake fiberoptic intubation for anticipated difficult airway.

\medterm{Difficult Intubation}-- Inability to perform direct laryngoscopy or endotracheal intubation despite appropriate expertise and equipment. Predictors include Mallampati Class III or IV, limited mouth opening (less than 3 cm), limited neck mobility, short thyromental distance (less than 6.5 cm), obesity, beard, facial trauma, and prior difficult intubation. The difficult airway algorithm emphasizes preoperative assessment, preparation, and sequential approach: preoxygenation, direct laryngoscopy with Macintosh blade, bougie, video laryngoscopy, supraglottic airway (intubating LMA, i-gel), and surgical airway (cricothyroidotomy). The cannot intubate, cannot oxygenate (CICO) situation requires immediate surgical airway access.

\medterm{Emergence Agitation}-- A state of confusion, restlessness, and inconsolable crying during emergence from anesthesia, occurring in 10 to 80 percent of pediatric patients and 5 to 10 percent of adults. Risk factors include sevoflurane induction, short procedures, parental separation, and preexisting behavioral disorders. Prevention includes gentle handling, parental presence during emergence, adequate analgesia, and minimizing stimulation. Treatment involves ensuring adequate oxygenation, adequate analgesia, and minimal stimulation. Treatment may require pharmacological intervention with propofol (0.5-1 mg/kg), midazolam (0.05 mg/kg), or dexmedetomidine (0.5 mcg/kg) in severe cases.

\medterm{Endotracheal Intubation}-- The definitive airway management procedure involving the insertion of a cuffed, flexible tube through the mouth (orotracheal) or nose (nasotracheal) past the vocal cords and into the trachea. It provides a secure, patent conduit for delivery of oxygen, inhalation anesthetics, and positive-pressure mechanical ventilation, while sealing the lower respiratory tract against aspiration of gastric contents and blood. Endotracheal tubes (ETTs) are sized by internal diameter (ID): typically 7.0 to 7.5 mm for adult females, and 7.5 to 8.5 mm for adult males; depth is typically taped at 21 cm at the incisors for females and 23 cm for males. Intubation is performed under direct laryngoscopy (Macintosh curved blade, Miller straight blade) or video laryngoscopy. Cuff pressure must be maintained at 20 to 30 cmH2O to prevent tracheal mucosal capillary ischemia (which occurs at $>30$ cmH2O) while preventing microaspiration. Definitive confirmation of correct tracheal placement requires continuous end-tidal CO2 detection with $\ge 6$ consistent waveforms, supported by direct visualization of the tube passing through the glottis, bilateral symmetric chest rise, and equal breath sounds. Complications include esophageal intubation, endobronchial (right mainstem) intubation, dental trauma, vocal cord injury, laryngospasm, and tracheal stenosis.

\medterm{Epidural Analgesia}-- A continuous regional analgesic technique that provides targeted segmental pain relief by infusing local anesthetics (e.g., bupivacaine 0.0625--0.125\% or ropivacaine 0.1--0.2\%) frequently combined with lipophilic opioids (fentanyl 2--4 mcg/mL or sufentanil) through a flexible indwelling catheter placed in the epidural space. Widely utilized for labor analgesia, major thoracic and upper abdominal surgery (thoracic epidural), and major lower extremity vascular or orthopedic procedures. Epidural analgesia provides superior dynamic pain relief compared to systemic opioids, attenuates the surgical neuroendocrine stress response, reduces postoperative pulmonary complications (atelectasis, pneumonia), promotes earlier return of bowel function, and decreases thromboembolic risk. Potential side effects include hypotension (resulting from preganglionic sympathetic vasomotor blockade), urinary retention, pruritus (opioid-mediated), and motor blockade. Catastrophic complications include epidural hematoma (dramatically increased in anticoagulated patients), epidural abscess, systemic local anesthetic toxicity from accidental intravascular migration, and total spinal blockade from accidental intrathecal migration. Contraindications include patient refusal, uncorrected coagulopathy or therapeutic anticoagulation, localized infection at the puncture site, and severe uncorrected hypovolemia.

\medterm{Epidural Anesthesia}-- A regional neuraxial anesthetic technique involving the injection of local anesthetics and adjuvants into the epidural space (between the ligamentum flavum and dura mater), producing reversible sensory, motor, and sympathetic blockade. Achieved by identifying the epidural space using the loss-of-resistance technique (with air or saline) via a Tuohy needle, followed by incremental injection of surgical doses of local anesthetic (e.g., bupivacaine 0.5\%, ropivacaine 0.5--0.75\%, or lidocaine 2\% with epinephrine) with or without catheter placement. Produces a gradual onset of blockade (15--25 minutes) compared to spinal anesthesia, with titratable level and duration. Side effects include sympathetic blockade-induced hypotension, motor blockade, and urinary retention. Complications include inadvertent dural puncture causing post-dural puncture headache (PDPH), systemic toxicity (LAST) from intravascular injection, epidural hematoma, and epidural abscess. Absolute contraindications include patient refusal, localized cutaneous infection, uncorrected coagulopathy, and severe hemodynamic shock.

\medterm{Epidural Blood Patch}-- A minimally invasive procedure used to treat post-dural puncture headache (PDPH) resulting from accidental dural puncture during epidural placement or spinal anesthesia. It involves injecting a small volume of the patient's autologous blood (typically 15-20 mL) into the epidural space at or near the site of the previous puncture. The blood clots over the dural hole, stopping further leakage of cerebrospinal fluid and rapidly resolving the headache.

\medterm{Etomidate}-- A carboxylated imidazole derivative used primarily as an intravenous hypnotic induction agent for general anesthesia, particularly in hemodynamically unstable, trauma, or critically ill patients. It acts by allosterically potentiating GABA-A receptor-mediated chloride currents in the central nervous system. Standard induction dose is 0.2 to 0.3 mg/kg IV, producing rapid unconsciousness within 30 to 60 seconds with a duration of 5 to 10 minutes due to redistribution. Its primary clinical advantage is remarkable cardiovascular stability: it preserves heart rate, systemic vascular resistance, mean arterial pressure, and cardiac output, causing minimal myocardial depression. It also reduces cerebral metabolic rate (CMRO2), cerebral blood flow, and intracranial pressure while maintaining cerebral perfusion pressure. Major limitations include: (1) Reversible, dose-dependent inhibition of $11\beta$-hydroxylase (the enzyme converting 11-deoxycortisol to cortisol in the adrenal cortex), causing adrenocortical suppression that persists for 6 to 24 hours even after a single induction dose; (2) High incidence of myoclonus during induction (occurring in up to 50--80 percent of unmedicated patients, mitigated by premedication with opioids or midazolam); (3) High rate of postoperative nausea and vomiting; and (4) Pain on injection due to formulation in propylene glycol.

\medterm{Extubation Criteria}-- Extubation requires the patient to meet specific criteria: fully awake and following commands, adequate tidal volume and respiratory rate, strong cough and gag reflex, adequate oxygenation (SpO2 greater than 95 percent on room air or baseline), hemodynamic stability, and return of neuromuscular function (TOF ratio greater than 0.9, sustained head lift for 5 seconds). Extubation should be performed at light anesthesia or fully awake. Awake extubation is preferred in difficult airway, obese, or high risk of aspiration (obstetric patients, obesity, gastro-oesophageal reflux disease, full stomach). Complications of extubation include laryngospasm, bronchospasm, aspiration, airway edema, and hypoxemia. Use of airway exchange catheters is recommended in patients with a known difficult airway to allow re-intubation if needed.

\medterm{Failed Intubation}-- Inability to successfully place an endotracheal tube after multiple attempts. In the operating room, a failed intubation algorithm is followed: call for help, maintain oxygenation (face mask, supraglottic airway), and decide to wake the patient (if non-emergent) or proceed with a supraglottic airway (if surgery is urgent). In the emergency department or ICU: difficult airway cart, video laryngoscopy, bougie, and surgical airway (cricothyroidotomy) if cannot intubate and cannot oxygenate. Failed intubation is more common in obstetric anesthesia (1 in 200-300) due to airway changes in pregnancy.

\medterm{Fasting Guidelines}-- Preoperative fasting guidelines minimize aspiration risk while reducing patient discomfort and catabolism. Clear liquids (water, clear juice without pulp, black coffee, sports drinks) are allowed up to 2 hours before surgery. Breast milk up to 4 hours. Infant formula up to 6 hours. Light meal (toast and clear liquid) up to 6 hours. Heavy/fatty meal up to 8 hours. The American Society of Anesthesiologists (ASA) fasting guidelines are evidence-based. Preoperative carbohydrate loading (clear carbohydrate drink) 2-3 hours before surgery reduces perioperative insulin resistance, thirst, and hunger without increasing aspiration risk, supporting enhanced recovery after surgery protocols.

\medterm{Fentanyl Derivatives}-- Opioid derivatives used in anesthesia include fentanyl (most common, 50 to 100 mcg IV), sufentanil (5 to 10 times more potent than fentanyl, 5 to 10 mcg IV, used in cardiac anesthesia), remifentanil (ultra-short-acting, metabolized by esterases, 0.1 to 0.5 mcg/kg/min infusion), alfentanil (shorter duration than fentanyl, used for brief procedures), and carfentanil (10,000 times more potent than fentanyl, used for wildlife immobilization). All opioids cause dose-dependent respiratory depression, histamine release, hypotension, and bradycardia at high doses and rapid administration. Fentanyl is lipid-soluble, allowing rapid brain penetration and making it suitable for transdermal delivery (fentanyl patch) for chronic pain. Fentanyl causes chest wall rigidity at rapid high-dose administration (greater than 150 mcg IV bolus), requiring positive pressure ventilation or neuromuscular blockade to overcome.

\medterm{General Anesthesia}-- A controlled, reversible state of unconsciousness, analgesia, amnesia, and muscle relaxation induced by pharmacologic agents, enabling surgical procedures without patient awareness or pain. The components of general anesthesia are induction (rapid onset of unconsciousness), maintenance (sustained anesthetic state), and emergence (recovery of consciousness). Induction agents include propofol (1.5 to 2.5 mg/kg IV, onset 15 to 30 seconds, duration 5 to 10 minutes, the most common induction agent), etomidate (0.2-0.3 mg/kg, haemodynamically neutral), and ketamine (1-2 mg/kg, dissociative anesthetic producing sympathetic stimulation and bronchodilation). Maintenance: volatile anesthetics (sevoflurane, isoflurane, desflurane) or total intravenous anesthesia (TIVA) with propofol and remifentanil infusions. Muscle relaxation can be achieved with depolarizing (succinylcholine) or non-depolarizing (rocuronium, vecuronium, atracurium) neuromuscular blocking agents. Emergence involves discontinuing anesthetics, reversing muscle relaxation, and extubating the patient once criteria are met.

\medterm{Hypertension in Anesthesia}-- Hypertension during anesthesia causes increased bleeding, myocardial ischemia, and stroke risk. Common causes include light anesthesia, pain, hypercarbia, hypoxia, bladder distension, shivering, and drug withdrawal (beta-blockers, clonidine). Treatment involves deepening anesthesia, providing analgesia (fentanyl, remifentanil), and antihypertensive agents (esmolol for tachycardia-associated hypertension, labetalol, nicardipine, nitroglycerin). In neurosurgery, blood pressure control prevents cerebral perfusion pressure management and avoidance of hypertension-induced bleeding.

\medterm{Hyperthermia in Anesthesia}-- Perioperative hyperthermia (temperature greater than 38 degrees Celsius) has multiple causes. Malignant hyperthermia (discussed separately) is the most urgent. Sepsis, thyroid storm, serotonin syndrome, neuroleptic malignant syndrome, and central fever are other causes. Fever increases metabolic rate by 10 to 13 percent per degree Celsius, increasing oxygen consumption and carbon dioxide production. Treatment involves identifying and addressing the underlying cause. Antipyretics (acetaminophen, and NSAIDs) are used. Cooling measures include forced air warming (paradoxically), cold intravenous fluids, ice packs, and cooling blankets. Avoid shivering if possible (meperidine 12.5-25 mg IV, dexmedetomidine). Identify and treat the underlying cause: dantrolene for malignant hyperthermia, beta-blockers for thyroid storm, cyproheptadine for serotonin syndrome.

\medterm{Hypotension in Anesthesia}-- Mean arterial pressure less than 65 mmHg or a decrease greater than 20 percent from baseline during anesthesia. Causes include anesthetic depth (volatile agents, propofol, opioids), hypovolemia (fasting, blood loss, third spacing), vasodilation (neuraxial blockade, sepsis, anaphylaxis), and cardiac dysfunction (myocardial depression, arrhythmias). Treatment involves identifying and addressing the cause: reducing anesthetic depth, fluid resuscitation, vasopressors (phenylephrine 50-100 mcg IV, ephedrine 5-10 mg IV, norepinephrine infusion), and inotropic support (epinephrine, dobutamine) if cardiac dysfunction is suspected or confirmed.

\medterm{Hypothermia in Anesthesia}-- A reduction in core body temperature below 36.0°C (96.8°F), occurring in 50 to 80 percent of surgical patients undergoing general or neuraxial anesthesia. General and regional anesthetics impair hypothalamic thermoregulatory control, widen the interthreshold range from 0.2°C to up to 4°C, and inhibit shivering and vasoconstriction. The drop in temperature occurs in three phases: Phase 1 (first hour), rapid core temperature decline (0.5--1.5°C) caused by core-to-peripheral heat redistribution from anesthetic-induced vasodilation; Phase 2 (subsequent 2--3 hours), slower linear decline due to radiant, convective, evaporative, and conductive heat loss exceeding metabolic heat production; and Phase 3, thermal plateau when vasoconstriction resets or heat loss equals production. Even mild hypothermia (34--36°C) triples the risk of surgical site infections (due to tissue hypoxia and impaired neutrophil function), induces coagulopathy (impaired platelet aggregation and clotting factor enzyme kinetics, increasing blood loss and transfusion requirements by ~20 percent), prolongs drug metabolism (delayed emergence), and triggers postoperative shivering, which increases oxygen consumption by 100 to 400 percent and elevates myocardial ischemia risk. Prevention utilizes active prewarming and intraoperative forced-air warming blankets, warmed intravenous fluids, and heat-and-moisture exchangers.

\medterm{Inhalational Anesthesia}-- The induction and maintenance of general anesthesia through the pulmonary administration of volatile liquids (sevoflurane, isoflurane, desflurane) or gases (nitrous oxide) delivered via an anesthesia breathing system. The depth of anesthesia depends directly on the partial pressure of the anesthetic gas in the brain ($P_{brain}$), which is in dynamic equilibrium with arterial blood ($P_{art}$) and alveolar gas ($P_{alv}$). Alveolar uptake is governed by three primary factors: (1) Blood-gas partition coefficient ($\lambda_{b/g}$), which reflects solubility; poorly soluble agents (desflurane 0.42, sevoflurane 0.65) achieve rapid rise in alveolar-to-inspired ratio ($F_A/F_I$), yielding fast induction and emergence, whereas soluble agents (isoflurane 1.4) dissolve readily in blood, delaying induction and emergence; (2) Alveolar ventilation, where increased minute ventilation accelerates the rise in $F_A/F_I$; and (3) Cardiac output, where increased pulmonary blood flow removes anesthetic from alveoli, paradoxically slowing the rise of $F_A/F_I$. Inhalation induction is standard in pediatric anesthesia using sweet-smelling, non-pungent sevoflurane, avoiding the distress of needle cannulation.

\medterm{Inotropes in Anesthesia}-- Inotropes increase myocardial contractility. Dobutamine is a beta-1 and beta-2 agonist increasing cardiac output with minimal effect on heart rate. Dose: 2 to 20 mcg/kg/min infusion. Milrinone is a phosphodiesterase-3 inhibitor increasing contractility and causing vasodilation. Dose: loading dose 50 mcg/kg over 10 minutes, then 0.375 to 0.75 mcg/kg/min. Epinephrine has potent beta-1 inotropic effects at low doses and alpha effects at high doses. Dose: 0.01 to 0.5 mcg/kg/min. Inotropes are used for cardiac dysfunction (myocardial depression, heart failure) and vasodilatory shock (sepsis, anaphylaxis). Adverse effects of inotropes include arrhythmias, increased myocardial oxygen consumption, and tolerance with prolonged use.

\medterm{Intraoperative Blood Transfusion}-- The intraoperative administration of allogeneic or autologous blood components to restore oxygen-carrying capacity and maintain intravascular volume during major hemorrhage. Transfusion thresholds are guided by physiological reserve: hemoglobin $<7$ g/dL in hemodynamically stable patients, or $<8$ g/dL in patients with pre-existing cardiovascular disease or active myocardial ischemia. One unit of packed red blood cells (PRBCs) increases hemoglobin by approximately 1 g/dL (hematocrit by 3 percent). In massive hemorrhage (>1 blood volume in 24 hours, or $>4$ units in 1 hour with ongoing bleeding), a Massive Transfusion Protocol (MTP) is activated, delivering balanced blood component therapy with PRBCs, fresh frozen plasma (FFP), and platelets in a 1:1:1 ratio. Point-of-care viscoelastic testing (thromboelastography [TEG] or rotational thromboelastometry [ROTEM]) guides targeted administration of cryoprecipitate (for fibrinogen $<150--200$ mg/dL), prothrombin complex concentrates, and antifibrinolytics. Complications include hypothermia, dilutional coagulopathy, hypocalcemia (citrate toxicity binding ionized calcium), hyperkalemia, transfusion-related acute lung injury (TRALI), and transfusion-associated circulatory overload (TACO).

\medterm{Intraoperative Neuromonitoring}-- Intraoperative neuromonitoring (IONM) uses electrophysiological techniques to monitor neural pathway integrity during surgery. Somatosensory evoked potentials (SSEP) monitor sensory pathways. Motor evoked potentials (MEP) monitor motor pathways. Electromyography (EMG) monitors cranial nerves and spinal roots. IONM is used in spine surgery (decompressive procedures, scoliosis correction), vascular surgery (carotid endarterectomy, aortic aneurysm repair), and skull base surgery. Changes in IONM signals require prompt communication with the surgical team to guide intervention (e.g., reducing retraction, altering surgical approach, or completing resection). anesthetic management must be tailored to IONM requirements, avoiding neuromuscular blocking agents during MEP monitoring and using total intravenous anesthesia (propofol and remifentanil) to minimise interference with evoked potential monitoring.

\medterm{Intraoperative Temperature Monitoring}-- The continuous measurement of core body temperature during surgical procedures exceeding 30 minutes to detect hypothermia, hyperthermia, and malignant hyperthermia. Reliable core measurement sites include the distal third of the esophagus (most accurate in intubated patients, positioned at the point of maximal heart sounds to avoid airway cooling artifact), the tympanic membrane (using thermocouple contact probes reflecting internal carotid artery temperature), the nasopharynx (placed in the posterior nasopharynx above the soft palate), the pulmonary artery (gold standard via Swan-Ganz catheter), and the urinary bladder (accurate in the absence of rapid cold diuresis). Peripheral sites (axilla, forehead, skin) underestimate core temperature by 1 to 2°C during vasoconstriction and are inaccurate for monitoring core changes. Continuous core monitoring is an essential ASA standard of care that guides active patient warming and temperature stabilization.

\medterm{Intravenous Anesthesia}-- The induction and maintenance of general anesthesia achieved through the administration of intravenous pharmacological agents, bypassing the pulmonary route. Intravenous induction agents (propofol, etomidate, ketamine, thiopental) produce rapid unconsciousness within one arm-to-brain circulation time (15--30 seconds) through high lipid solubility and high cerebral blood flow. Offset of action after a single bolus is mediated by rapid redistribution from high-perfusion vessel-rich organs (brain, heart) to intermediate muscle and poorly perfused fat tissue, rather than metabolic clearance. When anesthesia is maintained entirely with intravenous agents, it is termed Total Intravenous Anesthesia (TIVA), standardly combining propofol and remifentanil infusions delivered by target-controlled infusion (TCI) pumps. Intravenous anesthesia is particularly advantageous for neurosurgery, intraoperative neurophysiological monitoring (IONM, as TIVA does not abolish evoked potentials), history of malignant hyperthermia, and patients at high risk for postoperative nausea and vomiting.

\medterm{Isoflurane}-- A halogenated methyl ethyl ether volatile anesthetic agent that was long the global gold standard for general anesthesia maintenance. It has a minimum alveolar concentration (MAC) of 1.15\% in pure oxygen in young adults, making it more potent than sevoflurane and desflurane. Its blood-gas partition coefficient is 1.4, which confers moderate solubility and results in slower induction and emergence compared to sevoflurane and desflurane. Physiological features: (1) Neuroanesthesia: At concentrations $<1$ MAC, it preserves cerebral blood flow autoregulation while substantially reducing cerebral metabolic rate (CMRO2), producing an isoelectric EEG at 2 MAC, making it neuroprotective during periods of cerebral ischemia; (2) Cardiovascular: It produces dose-dependent peripheral vasodilation with preservation of cardiac output, though rapid concentration changes can cause modest tachycardia; in patients with critical coronary artery stenosis, high doses can theoretically induce 'coronary steal' by dilating healthy arterioles and diverting blood from underperfused ischemic beds; and (3) Respiratory: Potent bronchodilator, but possesses a mildly pungent odor that makes mask induction unpleasant compared to sevoflurane. It undergoes minimal hepatic metabolism (0.2\%) and triggers malignant hyperthermia in susceptible individuals.

\medterm{Ketamine}-- A phencyclidine (PCP) derivative that produces a state of 'dissociative anesthesia' characterized by profound analgesia, amnesia, catatonia, and dose-dependent hypnosis, while characteristically preserving protective airway reflexes, spontaneous respiration, and cardiovascular stability. Its primary mechanism of action is non-competitive antagonism of the N-methyl-D-aspartate (NMDA) receptor complex in the central nervous system, disrupting functional communication between the thalamocortical and limbic systems. Dosing: Intravenous induction dose is 1.0 to 2.0 mg/kg IV (onset 30--60 seconds, duration 10--15 minutes; or 4--6 mg/kg IM); subanesthetic analgesic infusion dose is 0.15 to 0.3 mg/kg IV bolus followed by 0.1 to 0.2 mg/kg/hr as a multimodal opioid-sparing adjunct. Hemodynamic effects: Unlike other hypnotics, ketamine stimulates the sympathetic nervous system via central sympathetic outflow and inhibition of neuronal catecholamine reuptake, resulting in increased heart rate, blood pressure, and cardiac output (making it the induction agent of choice in septic, hypovolemic, or hemorrhagic shock; however, it is a direct myocardial depressant unmasked in critically catecholamine-depleted patients). Respiratory effects: Potent bronchodilator (drug of choice for status asthmaticus and reactive airway disease), preserves functional residual capacity and pharyngeal reflexes, but increases tracheobronchial and salivary secretions (often co-administered with glycopyrrolate). Adverse effects include emergence delirium (vivid dreams, hallucinations, agitation occurring in 10--30 percent of adults, mitigated by co-administering midazolam or propofol), increased intraocular pressure, and transient increases in intracranial pressure.

\medterm{Laryngospasm}-- A protective reflex causing violent adduction of the vocal cords, resulting in complete airway obstruction. It is most common during light anesthesia, emergence, and in pediatric patients. Triggers include secretions, blood, airway manipulation, and suctioning. Signs include paradoxical chest-abdominal movement, suprasternal retractions, and absent air entry. Treatment includes positive pressure ventilation with 100 percent oxygen via face mask, jaw thrust, and deepening anesthesia with propofol (0.5-1 mg/kg IV) or succinylcholine (0.1-0.5 mg/kg IV). If oxygenation fails, cricothyroidotomy or tracheostomy is required. Prevention includes maintaining adequate depth of anesthesia before airway manipulation and avoiding airway stimulation during light planes of anesthesia.

\medterm{Lidocaine}-- An intermediate-acting amino-amide local anesthetic and Class Ib antiarrhythmic agent that reversibly blocks voltage-gated sodium channels in neuronal membranes, preventing action potential generation and impulse conduction. Extensively utilized in anesthesiology for: (1) Infiltration and minor nerve blocks (1--2\% solution; onset 2--5 min, duration 1--2 hours); (2) Airway topicalization (4\% solution or aerosol) prior to awake intubation; (3) Blunting the hemodynamic response to laryngoscopy and cough reflex during extubation (1.0--1.5 mg/kg IV 90 seconds prior); (4) Attenuating injection pain from propofol (20--40 mg IV); and (5) Perioperative intravenous systemic analgesia (bolus 1.5 mg/kg over 10 min followed by infusion 1.5--2.0 mg/kg/hr), which reduces postoperative pain, opioid requirements, ileus, and hospital length of stay in major abdominal and colorectal surgery. Maximum safe doses: 4.5 mg/kg plain (maximum 300 mg), or 7.0 mg/kg when combined with epinephrine (maximum 500 mg, where vasoconstriction limits systemic uptake). Metabolized in the liver by CYP1A2 and CYP3A4 to active metabolites (monoethylglycinexylidide). Signs of toxicity (LAST) progress from perioral numbness, metallic taste, and tinnitus to muscle twitching, grand mal seizures, respiratory arrest, and cardiovascular collapse.

\medterm{Lipid Emulsion Therapy}-- The definitive pharmacological antidote for local anesthetic systemic toxicity (LAST), as well as toxicity from other lipophilic drug overdoses (e.g., bupropion, calcium channel blockers, beta-blockers). The therapeutic mechanism involves the 'lipid sink' phenomenon (a lipid intravascular phase that captures and sequesters lipophilic local anesthetic molecules from target tissues in the brain and heart) and a metabolic effect (supplying abundant fatty acid substrates to reverse local-anesthetic-induced inhibition of cardiac mitochondrial oxidative phosphorylation). Standard resuscitation protocol utilizes a 20\% intravenous lipid emulsion (e.g., Intralipid): an initial IV bolus of 1.5 mL/kg (approx. 100 mL for a 70 kg adult) administered over 1 minute, followed immediately by a continuous infusion of 0.25 mL/kg/min (approx. 18 mL/min). If cardiovascular collapse persists, the bolus may be repeated once or twice every 3 to 5 minutes, and the infusion rate increased to 0.5 mL/kg/min, up to a recommended maximum total dose of 10--12 mL/kg over the first 30 minutes. Cardiopulmonary resuscitation must be continued throughout, and vasopressin, high-dose epinephrine (>1 mcg/kg), calcium channel blockers, and beta-blockers should be avoided. Propofol is formulated in a 10\% lipid emulsion but cannot be substituted for Intralipid because its hypnotic and vasodilatory properties will worsen cardiovascular collapse.

\medterm{Local Anesthetic}-- A medication that causes reversible loss of sensation in a specific, localized area of the body without affecting consciousness. Local anesthetics block voltage-gated sodium channels in nerve cell membranes, preventing action potential generation and propagation. The non-ionized (lipophilic) form crosses the nerve sheath; the ionized form binds to the sodium channel intracellularly, stabilizing it in the inactivated state. Small nerve fibers (A-delta and C, pain and temperature) are blocked before larger fibers (A-beta, touch and motor). Types: amides (lidocaine, bupivacaine, ropivacaine, prilocaine — metabolized in the liver) and esters (procaine, tetracaine, benzocaine — metabolized by plasma pseudocholinesterases). Maximum safe doses: lidocaine 4.5 mg/kg without epinephrine, 7 mg/kg with epinephrine; bupivacaine 2.25 mg/kg. Contraindications include allergy to same class, severe hepatic impairment (amides), and pseudocholinesterase deficiency (esters). Toxicity presents with CNS symptoms (perioral numbness, metallic taste, tinnitus, seizures, coma) followed by cardiovascular collapse (hypotension, arrhythmias, cardiac arrest). Treatment: stop injection, airway support, benzodiazepines for seizures, and Intralipid 20\% for cardiovascular toxicity. Addition of epinephrine prolongs duration and reduces systemic absorption.

\medterm{Local Anesthetic Systemic Toxicity}-- A life-threatening complication caused by elevated plasma concentration of local anesthetic, most commonly from inadvertent intravascular injection or rapid absorption. Presents with progressive CNS symptoms (perioral numbness, metallic taste, tinnitus, agitation, seizures) followed by cardiovascular effects (arrhythmias, bradycardia, hypotension, cardiac arrest). Cardiovascular collapse typically precedes CNS symptoms with bupivacaine (highly cardiotoxic). Prevention: test dose, incremental injection, aspiration, and ultrasound guidance. Treatment: STOP injection, call for help, airway management, lipid emulsion therapy (Intralipid 20\%: 1.5 mL/kg bolus, then 0.25 mL/kg/min infusion). Avoid vasopressin and calcium channel blockers. Continue CPR until lipid levels are therapeutic.

\medterm{Malignant Hyperthermia}-- A rare, life-threatening pharmacogenetic skeletal muscle disorder inherited in an autosomal dominant pattern with variable penetrance, triggered by volatile inhalational anesthetics (sevoflurane, desflurane, isoflurane, halothane) and the depolarizing muscle relaxant succinylcholine. The underlying defect involves mutations in the ryanodine receptor gene (RYR1 in >70 percent of cases) or the CACNA1S gene on the sarcoplasmic reticulum, resulting in uncontrolled calcium efflux into the myoplasm, triggering persistent, hypermetabolic skeletal muscle contraction, excessive glycogenolysis, and profound heat production. Clinical manifestations: The earliest and most sensitive sign in a ventilated patient is an abrupt, unexplained, exponential rise in end-tidal CO2 (refractory to increased ventilation), accompanied by sinus tachycardia, tachypnea, masseter muscle rigidity (masseter spasm following succinylcholine), skin mottling, and profuse diaphoresis. Late signs include rapidly escalating hyperthermia (rising 1--2°C every 5 minutes, exceeding 41°C), mixed respiratory and metabolic acidosis with severe hyperkalemia, rhabdomyolysis (hyperkalemia, elevated creatine kinase, myoglobinuria, acute renal failure), and disseminated intravascular coagulation (DIC). Management protocol: (1) Discontinue all triggering agents immediately; (2) Hyperventilate with 100\% oxygen at maximum fresh gas flow ($\ge 10$ L/min) using activated charcoal filters or a clean circuit; (3) Administer dantrolene sodium immediately at an initial bolus of 2.5 mg/kg IV, repeated every 5--10 minutes up to 10 mg/kg until physiological parameters normalize (dantrolene acts by directly blocking the ryanodine receptor, arresting calcium efflux); (4) Treat hyperkalemia aggressively (calcium chloride/gluconate, insulin with dextrose, sodium bicarbonate); (5) Initiate active external and internal core cooling (cold IV saline, ice packs to axillae/groin, bladder/gastric lavage), terminating cooling when temperature reaches 38°C to prevent hypothermia; (6) Maintain urine output $>1--2$ mL/kg/hr with fluids and mannitol/furosemide to prevent myoglobinuric renal tubular necrosis. ICU monitoring for 24--48 hours is mandatory due to a 20--25\% recrudescence rate.

\medterm{Mallampati Classification}-- A simple scoring system used to predict the ease of endotracheal intubation based on the anatomical structures visible in the oral cavity when the patient opens their mouth wide and protrudes their tongue without phonating. Class I: soft palate, fauces, uvula, and tonsillar pillars visible. Class II: soft palate, fauces, and uvula visible. Class III: soft palate and base of uvula visible. Class IV: only hard palate visible. Higher Mallampati classes (III and IV) are associated with a higher likelihood of a difficult airway.

\medterm{Mechanical Ventilation}-- The delivery of artificial respiratory support through positive-pressure breaths via an endotracheal tube or tracheostomy to maintain alveolar ventilation and arterial oxygenation. Standard ventilator modes include: Volume-Controlled Ventilation (VCV, delivering a preset tidal volume with variable airway pressure depending on compliance and resistance); Pressure-Controlled Ventilation (PCV, delivering breaths up to a preset inspiratory pressure with variable tidal volume, reducing barotrauma risk); Pressure Support Ventilation (PSV, a spontaneous mode where the patient triggers the breath and the ventilator assists with a set pressure boost); and Synchronized Intermittent Mandatory Ventilation (SIMV, delivering mandatory breaths synchronized with spontaneous efforts). Evidence-based lung-protective ventilation in the operating room and intensive care unit utilizes low tidal volumes (6--8 mL/kg of predicted ideal body weight, rather than actual weight), positive end-expiratory pressure (PEEP 5--10 cmH2O to prevent atelectrauma), recruitment maneuvers, and limitation of plateau airway pressure to $<30$ cmH2O and driving pressure ($\Delta P = P_{plat} - PEEP$) to $<15$ cmH2O to minimize ventilator-induced lung injury (VILI).

\medterm{Mendelson Syndrome}-- A historical eponym for severe chemical aspiration pneumonitis, first described by obstetrician Curtis Mendelson in 1946 in pregnant patients who aspirated acidic gastric contents during ether anesthesia. Classically defined by the acute pulmonary aspiration of more than 25 mL ($>0.4$ mL/kg) of gastric juice with a pH below 2.5, triggering severe chemical injury, bronchospasm, and non-cardiogenic pulmonary edema. See \hyperlink{Aspiration Pneumonitis}{Aspiration Pneumonitis}.

\medterm{Minimum Alveolar Concentration (MAC)}-- The alveolar concentration of an inhalational anesthetic at 1 atmosphere of pressure that prevents movement in 50\% of patients in response to a noxious surgical stimulus. It serves as the standard measure of potency for volatile anesthetics (e.g., sevoflurane has a MAC of ~2.0\%, while desflurane has a MAC of ~6.0\% in adults). MAC values are additive and are decreased by factors such as advanced age, pregnancy, hypothermia, and co-administration of opioids or sedatives.

\medterm{Monitored Anesthesia Care}-- A distinct service where an anesthesia professional provides premedication, administration of analgesics and/or sedatives, continuous monitoring of the patient's cardiorespiratory status, and intervention for any emergent or unavoidable complications that may arise during a diagnostic or therapeutic procedure. MAC is used for endoscopy, colonoscopy, interventional radiology, and minor surgical procedures. The patient may progress from moderate sedation to deep sedation, where airway intervention may be required. MAC requires at least one qualified anesthesia professional continuously present to monitor vital signs, administer sedation and analgesia, and be prepared to manage the airway and haemodynamic emergencies at any time during the procedure.

\medterm{Multimodal Analgesia}-- The systematic integration of multiple pharmacological agents with distinct mechanisms of action and/or regional analgesic techniques to target different neuroanatomical sites along the pain pathway (transduction, transmission, modulation, perception). The primary objective is to achieve additive or synergistic analgesia while minimizing reliance on systemic opioids, thereby significantly reducing opioid-related adverse effects (hypoventilation, sedation, postoperative nausea and vomiting, ileus, urinary retention, tolerance, and physical dependence). Core systemic components in modern Enhanced Recovery After Surgery (ERAS) pathways combine: around-the-clock non-opioid baseline analgesics (acetaminophen, NSAIDs or selective COX-2 inhibitors), membrane-stabilizing and neuromodulatory adjuvants (gabapentinoids, intravenous lidocaine infusions, low-dose ketamine infusions, dexamethasone, and alpha-2 agonists such as dexmedetomidine), and site-specific regional interventions (wound infiltration, fascial plane blocks, peripheral nerve blocks, or neuraxial blockade). Systemic opioids are reserved for breakthrough pain only.

\medterm{Nasopharyngeal Airway}-- A soft, flexible, hollow rubber or silicone tube (nasal trumpet) designed to be inserted through the external nares along the floor of the nasal cavity into the posterior pharynx to establish and maintain upper airway patency. Sizing is determined by measuring the distance from the tip of the patient's nose to the tragus of the ear (or angle of the mandible), with the internal diameter roughly approximating the patient's fifth (little) digit. Prior to insertion, the airway is thoroughly lubricated with water-soluble surgical gel (or topical vasoconstrictor/local anesthetic like lidocaine-phenylephrine to minimize epistaxis) and introduced with the bevel facing the nasal septum, advancing gently along the floor of the nose perpendicular to the face (not directed cephalad). Unlike an oropharyngeal airway, a nasopharyngeal airway is much better tolerated in patients who are semi-conscious, have an active gag reflex, or suffer from severe trismus, clenched jaw, or oral trauma. Contraindications include suspected basilar skull fractures (characterized by CSF rhinorrhea, Battle sign, or raccoon eyes, where intracranial tube entry has been reported), severe nasal trauma, and severe coagulopathy.

\medterm{Neostigmine}-- A synthetic quaternary ammonium compound that functions as a reversible acetylcholinesterase inhibitor (anticholinesterase), utilized as the traditional pharmacological agent for reversing non-depolarizing neuromuscular blockade. By competitively binding to the active esteratic site of the acetylcholinesterase enzyme, it prevents the hydrolysis of acetylcholine within the neuromuscular synaptic cleft, causing local acetylcholine accumulation that competes with and displaces non-depolarizing relaxants (e.g., rocuronium, vecuronium, cisatracurium) from nicotinic receptors. Clinical dosing: 0.04 to 0.07 mg/kg IV (maximum recommended dose 5 mg), with an onset of 5 to 10 minutes and peak effect at 10 minutes. It exhibits a pharmacodynamic 'ceiling effect': once acetylcholinesterase is 100\% inhibited, additional drug provides no further reversal and can paradoxically induce depolarizing weakness; consequently, it should only be administered when spontaneous recovery has already commenced (minimum 2, preferably 4 twitches on train-of-four monitoring). Because neostigmine causes systemic acetylcholine accumulation at muscarinic receptors, it triggers profound cholinergic side effects (bradycardia, bronchospasm, hypersalivation, miosis, increased bowel motility); therefore, it must always be co-administered with a muscarinic anticholinergic agent, standardly glycopyrrolate (0.2 mg per 1 mg neostigmine) or atropine, matching the onset of action.

\medterm{Neuromuscular Blocking Agents (NMBAs)}-- Medications administered intraoperatively to induce skeletal muscle relaxation and facilitate endotracheal intubation and surgical exposure. They act at the nicotinic acetylcholine receptors of the neuromuscular junction and are classified as depolarizing (e.g., succinylcholine, which acts as an agonist and causes transient fasciculations) or non-depolarizing (e.g., rocuronium, vecuronium, cisatracurium, which act as competitive antagonists). Their effect is monitored using a train-of-four (TOF) stimulator and can be reversed with anticholinesterases (e.g., neostigmine) or selective relaxant binding agents (e.g., sugammadex).

\medterm{Neuromuscular Junction}-- The specialized cholinergic chemical synapse between the unmyelinated terminal axon of an alpha motor neuron and the motor endplate of a skeletal muscle fiber. Action potential arrival triggers voltage-gated calcium influx into the presynaptic terminal, releasing acetylcholine (ACh) quanta from synaptic vesicles into the 20--50 nm synaptic cleft. ACh binds to postsynaptic nicotinic acetylcholine receptors (pentameric ligand-gated ion channels composed of two $\alpha$, one $\beta$, one $\delta$, and one $\epsilon$ subunit in adult muscle), inducing a conformational change that permits sodium and calcium influx and potassium efflux. The resulting endplate potential (EPP) depolarizes the sarcolemma, activating voltage-gated sodium channels and triggering muscle contraction. Acetylcholinesterase (AChE) anchored to the synaptic basal lamina rapidly hydrolyzes ACh into acetate and choline within milliseconds, terminating receptor activation. Neuromuscular blocking agents act at this site: succinylcholine binds as an agonist causing sustained Phase I depolarization, while non-depolarizing agents (rocuronium, vecuronium, cisatracurium) competitively antagonize ACh binding, preventing depolarization and causing flaccid paralysis.

\medterm{Nitrous Oxide}-- An inorganic, non-flammable, colorless gas ($N_2O$) with a pleasant, sweet odor, utilized as an adjuvant inhalational anesthetic and analgesic agent. It exerts moderate analgesia primarily through NMDA receptor inhibition and activation of descending noradrenergic pain pathways, but has remarkably low anesthetic potency, with a minimum alveolar concentration (MAC) of 104\%, meaning it cannot produce surgical general anesthesia alone under normal atmospheric conditions. It has an exceptionally low blood-gas partition coefficient (0.47), conferring extremely rapid induction and emergence. Clinical principles: (1) Second Gas Effect: Administration of high concentrations of N2O (e.g., 70\%) accelerates the rate of alveolar rise of a co-administered volatile anesthetic (e.g., sevoflurane); (2) Closed Gas Space Expansion: N2O is 34 times more soluble in blood than nitrogen; it diffuses into air-filled, compliant, or non-compliant closed body cavities faster than nitrogen can diffuse out, causing rapid expansion or severe pressure elevation. Consequently, N2O is strictly contraindicated in pneumothorax, bowel obstruction, air embolism, middle ear surgery, intraocular gas bubbles (SF6/C3F8), and pneumocephalus; (3) Diffusion Hypoxia: Following N2O discontinuation, large volumes rapidly pour from blood into alveoli, diluting alveolar oxygen and CO2; patients must receive 100\% oxygen for 5--10 minutes at emergence to prevent hypoxia; and (4) Methionine Synthase Inactivation: Irreversibly oxidizes the cobalt atom of vitamin B12, inhibiting methionine synthase and folate metabolism, predisposing to megaloblastic anemia, neurotoxicity (subacute combined degeneration of the spinal cord), and elevated homocysteine with prolonged or repeated exposure.

\medterm{Nonsteroidal Anti-inflammatory Drugs (NSAIDs)}-- A major class of non-opioid analgesic, antipyretic, and anti-inflammatory medications widely utilized as a cornerstone of perioperative multimodal analgesia. They act by inhibiting cyclooxygenase enzymes (COX-1 and COX-2), preventing the conversion of arachidonic acid to algesic prostaglandins (particularly PGE2 and PGI2), prostacyclin, and thromboxane. Classified into: (1) Non-selective COX inhibitors (e.g., ketorolac 15--30 mg IV q6h, ibuprofen 400--800 mg IV), which inhibit both COX-1 (constitutive, gastroprotective, and platelet-activating) and COX-2 (inducible in inflammation); and (2) COX-2 selective inhibitors (coxibs, e.g., celecoxib oral, parecoxib IV), which preserve platelet thromboxane production and reduce gastrointestinal mucosal toxicity. NSAIDs provide potent somatic and inflammatory analgesia and reduce postoperative opioid consumption by 20 to 40 percent. Adverse effects and contraindications: (1) Renal dysfunction: inhibition of vasodilatory prostaglandins blunts renal autoregulation, risking acute kidney injury in hypovolemia, dehydration, or pre-existing renal disease; (2) Platelet inhibition: non-selective agents reversibly inhibit COX-1-dependent thromboxane A2, prolonging bleeding time (contraindicated in active surgical hemorrhage, coagulopathy, or intracranial surgery); (3) Peptic ulceration and gastritis; (4) Bronchospasm in aspirin-exacerbated respiratory disease (AERD); and (5) Potential inhibition of osteogenesis, cautioning against prolonged high-dose use in spinal fusion or non-union fractures.

\medterm{Opioid Reversal}-- The pharmacological antagonism of opioid-induced central nervous system and respiratory depression, achieved primarily with naloxone, a pure competitive antagonist at mu, kappa, and delta opioid receptors. In postoperative or emergency settings, naloxone is administered intravenously in titrated increments (0.04 to 0.1 mg IV every 2--3 minutes) to restore adequate spontaneous ventilation and protective airway reflexes without precipitating acute unbearable surgical pain or sympathetic surge. The terminal elimination half-life of naloxone is short (30--80 minutes), which is substantially shorter than that of most opioids (e.g., morphine, hydromorphone, methadone, fentanyl patches); consequently, patients must be monitored continuously for renarcotization and recurrent hypoventilation, which may necessitate repeat boluses or a continuous IV infusion (typically calculated at two-thirds of the initial successful reversal bolus per hour). Longer-acting oral antagonists include naltrexone (half-life ~4 hours, active metabolite 6-beta-naltrexol ~13 hours) and nalmefene (half-life 8--11 hours). Caution: Rapid, high-dose opioid reversal in opioid-dependent patients can precipitate acute withdrawal, severe hypertension, tachycardia, cardiac arrhythmias, and non-cardiogenic pulmonary edema.

\medterm{Oropharyngeal Airway}-- A curved, rigid plastic or rubber medical device inserted through the oral cavity to displace the tongue and epiglottis anteriorly away from the posterior pharyngeal wall, maintaining a patent upper airway in an unconscious patient. Common designs include the Guedel airway (tubular with a central lumen) and Berman airway (channeled along the sides). Proper sizing is essential: measured from the corner of the patient's mouth to the angle of the mandible (or tragus of the ear); an undersized airway pushes the tongue posteriorly, worsening obstruction, while an oversized airway impinges on the epiglottis, potentially triggering laryngospasm or vocal cord damage. Insertion in adults is performed by inserting the airway upside-down (tip directed toward the hard palate) and rotating it 180 degrees as it passes past the soft palate, or directly using a tongue depressor to avoid palatal trauma. Crucially, it is strictly contraindicated in conscious or semi-conscious patients with an intact gag reflex, as it reliably triggers retching, vomiting, aspiration, or laryngospasm. It also functions as a bite block to prevent patient occlusion or damage to an endotracheal tube.

\medterm{Other Pain Conditions}-- Other pain types encompass pain conditions not primarily classified as nociceptive (somatic or visceral) or classic neuropathic pain, including nociplastic pain (altered nociception without clear tissue or nerve damage, such as fibromyalgia, chronic tension-type headache, temporomandibular disorders, and chronic pelvic pain), complex regional pain syndrome (CRPS Type I and II—persistent regional pain disproportionate to inciting event with autonomic and motor features), central pain syndromes (post-stroke pain, spinal cord injury pain), and chronic visceral pain syndromes (interstitial cystitis/bladder pain syndrome, irritable bowel syndrome, chronic pancreatitis, endometriosis). These conditions share mechanisms of central sensitization, dysfunctional descending modulation, and psychologic contributors. Diagnosis is clinical, often requiring exclusion of alternative pathology. Management is multimodal: patient education, physical therapy, cognitive-behavioral therapy, medications (tricyclic antidepressants, gabapentinoids, SNRIs), and neuromodulation; opioids are generally ineffective and not recommended for chronic non-cancer pain.

\medterm{Pain}-- An unpleasant sensory and emotional experience associated with actual or potential tissue damage, serving as a protective alerting system but becoming maladaptive when chronic (lasting >3 months). The biopsychosocial model recognizes nociception (tissue damage), pain perception (sensory-discriminative, affective-motivational, cognitive-evaluative dimensions), and pain behaviors influenced by genetics, psychology, social context, and comorbidities. Classification: nociceptive (somatic—well-localized, aching; visceral—diffuse, cramping), neuropathic (nerve injury—burning, shooting, allodynia), and nociplastic (altered nociception without clear lesion—fibromyalgia). Acute pain resolves with healing; chronic pain affects ~20\% of adults and is a leading cause of disability. Assessment: PQRST (provoking/palliating, quality, region/radiation, severity, timing), functional impact, and psychologic screening (depression, catastrophizing, kinesiophobia). Management requires an individualized multimodal approach.

\medterm{Pain Assessment Scales}-- Pain assessment uses standardized scales to quantify pain intensity. Numeric rating scale (NRS) asks patients to rate pain from 0 (no pain) to 10 (worst pain imaginable). Visual analog scale (VAS) is a 10 cm line with anchors. Faces pain scale (Wong-Baker FACES) uses facial expressions for children and cognitively impaired patients. FLACC scale (Face, Legs, Activity, Cry, Consolability) assesses pain in non-verbal children. Critical-care Pain Observation Tool (CPOT) assesses pain in mechanicallyventilated patients in the intensive care unit. PAINAD (Pain Assessment in Advanced Dementia) assesses pain in elderly patients with dementia. Pain should be assessed at rest and with movement. Standardized scales improve pain documentation, treatment consistency, and outcomes. The goal is to maintain pain scores below 3-4 out of 10, allowing the patient to function and participate in recovery.

\medterm{Pain Relief (Analgesia)}-- Achieved through pharmacologic, interventional, physical, and psychologic modalities targeting nociceptive transmission and pain perception. Pharmacologic: World Health Organization (WHO) analgesic ladder recommends non-opioids (acetaminophen, NSAIDs) for mild pain, weak opioids (tramadol, codeine) for moderate pain, and strong opioids (morphine, oxycodone, hydromorphone, fentanyl) for severe pain, with adjuvant medications (antidepressants, anticonvulsants, local anesthetics) for neuropathic pain. Interventional procedures: nerve blocks (facet, epidural, peripheral nerve), radiofrequency ablation, spinal cord stimulation, intrathecal pumps. Physical modalities: ice/heat, transcutaneous electrical nerve stimulation (TENS), physical therapy, massage, acupuncture. Psychologic: cognitive-behavioral therapy, biofeedback, relaxation techniques, mindfulness meditation. Opioid prescribing requires careful risk-benefit assessment given addiction potential (especially with chronic use), respiratory depression, tolerance, and opioid-induced hyperalgesia. Multimodal analgesia—combining agents with different mechanisms—improves efficacy and reduces opioid requirements (opioid-sparing effect).

\medterm{Patient-Controlled Analgesia}-- Patient-controlled analgesia (PCA) allows patients to self-administer small doses of analgesic medication within preset safety limits. IV PCA typically uses morphine (demand 1 mg, lockout 6 to 10 minutes, hourly limit 4 mg), hydromorphone (demand 0.2 mg, lockout 6 to 10 minutes), or fentanyl (demand 10 to 20 mcg, lockout 6 to 10 minutes). PCA provides more consistent pain control than nurse-administered boluses and gives patients a sense of control. Basal infusion is controversial and increases the risk of respiratory depression, especially in opioid-naive patients and those with obstructive sleep apnea. PCA by proxy (family members or nurses activating the button) is dangerous and should be avoided. Contraindications include patient inability to understand the device (cognitive impairment, age <5 years), allergy, and opioid tolerance. Epidural PCA (PCEA) combines a background infusion with patient-controlled boluses and is used for postoperative and labour analgesia.

\medterm{Pediatric Anesthesia}-- Pediatric anesthesia requires age-appropriate dosing, equipment, and techniques. Children have higher metabolic rates, increased oxygen consumption, smaller functional residual capacity, and higher risk of desaturation. Airway anatomy differs (larger tongue, more anterior larynx, narrower subglottic area). Induction is commonly performed via inhalation (sevoflurane in oxygen) or intravenous. Dosing is based on weight (mg/kg) with adjustments for neonates and infants. Temperature management is critical in neonates and infants due to high surface area-to-volume ratio. Emergence delirium is more common in children (sevoflurane) and is managed with propofol, dexmedetomidine, or gentle awakening. Postoperative nausea and vomiting is also common and treated with ondansetron, dexamethasone, and total intravenous anesthesia.

\medterm{Peripheral Nerve Blocks}-- Peripheral nerve blocks provide targeted anesthesia and analgesia for specific surgical sites. Ultrasound-guided blocks improve success rates and reduce complications. Common blocks include interscalene (shoulder surgery), supraclavicular (arm surgery), infraclavicular (arm and hand), axillary (hand surgery), femoral (knee surgery), adductor canal (knee surgery with preserved quadriceps function), sciatic (foot and ankle), popliteal (foot and ankle), and ilioinguinal/iliohypogastric (inguinal hernia repair). Local anesthetics (bupivacaine 0.25-0.5\%, ropivacaine 0.2-0.5\%, lidocaine 1-2\%) are used with or without epinephrine to prolong block duration and reduce systemic absorption. Ultrasound guidance has become the standard of care, reducing the risk of inadvertent vascular puncture, intraneural injection, and pneumothorax while improving block success rates and reducing local anesthetic volume requirements.

\medterm{Postanesthesia Care Unit}-- A specialized, critical care recovery environment immediately adjacent to the operating suite where patients are monitored while recovering from the physiological effects of general, regional, or monitored anesthesia care. Staffed by specialized perianesthesia nurses with 1:1 or 1:2 nurse-to-patient ratios, the PACU provides continuous monitoring of airway patency, pulse oximetry, electrocardiography, noninvasive blood pressure, and core temperature. Clinical assessment on admission includes hemodynamic stability, level of consciousness, surgical wound and drain status, fluid balance, and pain control. Recovery and readiness for discharge are objectively evaluated using scoring systems, most commonly the Modified Aldrete Score, which grades 5 domains (0, 1, or 2 points each): (1) Activity (moves 4, 2, or 0 extremities voluntarily); (2) Respiration (deep breath/cough, dyspnea/shallow, or apnea); (3) Circulation (blood pressure within 20\%, 20--50\%, or $>50\%$ of baseline); (4) Consciousness (fully awake, arousable on calling, or unresponsive); and (5) Oxygen Saturation (SpO2 $>92\%$ on room air, requires oxygen to maintain $>90\%$, or $<90\%$ despite oxygen). A score of $\ge 9$ (or return to preoperative baseline) is generally required for transfer to an inpatient surgical floor. Common PACU emergencies include upper airway obstruction (managed with jaw thrust, chin lift, oral/nasal airways), postoperative nausea and vomiting (treated with ondansetron, droperidol), hypothermia, acute surgical hemorrhage, and severe pain (treated with PCA or regional analgesia).

\medterm{Post-Dural Puncture Headache}-- A positional headache occurring after accidental dural puncture during epidural placement or lumbar puncture, caused by cerebrospinal fluid leakage and intracranial hypotension. Symptoms include frontal or occipital headache worse when upright and relieved when supine, neck stiffness, tinnitus, photophobia, and nausea. Incidence is 1 to 2 percent with 25 gauge pencil-point needles, 5 to 10 percent with 22 gauge cutting needles. Treatment includes conservative: bed rest, hydration, caffeine (500 mg IV or oral), and analgesics. Invasive: epidural blood patch (15-20 mL of autologous blood injected into the epidural space at the level of the dural puncture), which is highly effective (80-95\% success rate). Prevention: use of pencil-point (Whitacre, Sprotte) instead of cutting needles (Quincke), small gauge (25-27 G), and orientation of the needle bevel parallel to the dural fibers.

\medterm{Postoperative Cognitive Dysfunction}-- Decline in cognitive function (memory, concentration, executive function) after surgery and anesthesia, more common in elderly patients. POCD can be early (weeks to months) or persistent (>1 year). Incidence: 25-40\% in elderly patients at hospital discharge, 10-15\% at 3 months. Risk factors: age, lower education level, preoperative cognitive impairment, prolonged surgery, respiratory complications, and postoperative delirium. Pathophysiology: neuroinflammation (surgical stress response, cytokines), anesthetics (GABAergic drugs), and microemboli. Prevention: minimize depth and duration of anesthesia, avoid anticholinergic drugs, early mobilization, and cognitive stimulation. Management: neuropsychological evaluation and cognitive rehabilitation.

\medterm{Postoperative Delirium}-- An acute, fluctuating disturbance of attention, awareness, and cognition occurring in the immediate postoperative period (typically peaking between postoperative days 1 and 3), particularly prevalent in elderly surgical patients (15--50 percent incidence). It is clinically categorized into three motor subtypes: hyperactive (agitation, restlessness, hallucinations, combativeness; ~25\%), hypoactive (lethargy, apathy, psychomotor slowing, decreased responsiveness; ~50\%, the most common and under-recognized form), and mixed (~25\%). Diagnosed using validated bedside tools such as the Confusion Assessment Method (CAM or CAM-ICU). Risk factors include advanced age ($>65$ years), baseline cognitive impairment or dementia, sensory deficits, polypharmacy, high-risk surgical procedures (orthopedic hip fracture, cardiac surgery), and exposure to deliriogenic drugs (anticholinergics, benzodiazepines, meperidine). Delirium significantly increases hospital length of stay, postoperative complications, institutionalization rates, long-term cognitive decline that may persist for months to years, and mortality. Prevention relies on non-pharmacological multicomponent protocols: frequent reorientation, family presence, early mobilization, sleep-wake cycle preservation, sensory aids (glasses, hearing aids), and avoidance of deliriogenic medications.

\medterm{Postoperative Nausea and Vomiting}-- Nausea, retching, or vomiting occurring within 24-48 hours after surgery. PONV is one of the most common and distressing complications. Risk factors (Apfel score): female gender, non-smoker, history of motion sickness/PONV, and postoperative opioids. 0-1 risk factors: 10-20\% incidence; 2: 40\%; 3: 60\%; 4: 80\%. Prophylaxis: multimodal approach — 5-HT3 antagonists (ondansetron 4 mg), dexamethasone 4-8 mg, droperidol 0.625 mg, scopolamine patch, and total intravenous anesthesia with propofol. Treatment: 5-HT3 antagonist + dexamethasone + droperidol for rescue; consider NK1 antagonists (aprepitant) for high-risk patients.

\medterm{Postoperative Pain Management}-- The systematic relief of acute post-surgical pain using an evidence-based multimodal approach that combines diverse pharmacological agents and regional analgesic techniques to optimize comfort, facilitate early mobilization, and minimize opioid-related morbidity. Core systemic components include around-the-clock non-opioids: acetaminophen (1 g IV/oral q6h) and nonsteroidal anti-inflammatory drugs (NSAIDs, e.g., ketorolac 15--30 mg IV q6h or celecoxib) to provide foundational anti-inflammatory analgesia. Adjuvant agents include gabapentinoids (pregabalin, gabapentin) for neuropathic/hyperalgesic components, low-dose intravenous ketamine infusions (0.1--0.2 mg/kg/hr) to inhibit NMDA receptor-mediated central sensitization, and dexamethasone (4--8 mg IV) for anti-inflammatory and antiemetic synergy. Regional analgesia—including continuous epidural infusions, peripheral nerve blocks (e.g., femoral, adductor canal, sciatic, interscalene), and fascial plane blocks (TAP, ESP, quadratus lumborum blocks)—provides profound, site-specific, targeted analgesia. Intravenous patient-controlled analgesia (PCA) with opioids (morphine, hydromorphone, fentanyl) is reserved for severe breakthrough pain. Effective postoperative pain control reduces pulmonary complications, blunts the surgical stress response, and prevents transition to chronic post-surgical pain.

\medterm{Post-Surgical Pain}-- Acute nociceptive, neuropathic, and inflammatory pain resulting directly from surgical tissue trauma, incision, visceral traction, dissection, and ischemia. Tissue injury triggers the local release of an 'inflammatory soup' of algesic mediators (prostaglandins, bradykinin, histamine, cytokines, nerve growth factor, substance P) that sensitize peripheral high-threshold primary afferent nociceptors (peripheral sensitization). Repeated nociceptive bombardment of dorsal horn neurons in the spinal cord induces central sensitization (the 'wind-up' phenomenon), resulting in primary and secondary hyperalgesia and allodynia. Uncontrolled post-surgical pain induces autonomic stress responses (hypertension, tachycardia, myocardial ischemia), limits respiratory excursion (promoting atelectasis, hypoxemia, and pneumonia), delays ambulation (increasing deep vein thrombosis risk), and dramatically elevates the incidence of chronic post-surgical pain (CPSP, occurring in 10 to 50 percent of patients undergoing thoracotomy, mastectomy, amputation, or inguinal hernia repair). Optimal management follows Enhanced Recovery After Surgery (ERAS) principles, employing preventive and multimodal analgesia combining regional blockade, around-the-clock non-opioids, and targeted rescue analgesia.

\medterm{Preoperative Assessment}-- A comprehensive clinical evaluation performed prior to anesthesia and surgery to identify patient comorbidities, stratify perioperative risk, formulate an individualized anesthetic plan, and optimize medical conditions. It encompasses a detailed medical history (cardiorespiratory disease, functional exercise capacity measured in metabolic equivalents [METs], prior anesthesia history including difficult airway, malignant hyperthermia susceptibility, and family history), a targeted physical examination (vital signs, cardiopulmonary auscultation, body mass index, and comprehensive airway evaluation), and selective, indication-based preoperative laboratory investigations. Risk assessment utilizes validated scoring systems: (1) Revised Cardiac Risk Index (Lee RCRI) for major adverse cardiac events; (2) STOP-BANG questionnaire for obstructive sleep apnea risk; and (3) the ASA Physical Status Classification System, grading fitness from ASA I (normal healthy patient), ASA II (mild systemic disease), ASA III (severe systemic disease), ASA IV (severe systemic disease that is a constant threat to life), ASA V (moribund patient not expected to survive 24 hours without surgery), and ASA VI (declared brain-dead organ donor), with the suffix 'E' appended for emergency procedures. Preoperative fasting instructions follow evidence-based guidelines (2 hours for clear liquids, 6 hours for light meals, 8 hours for fatty meals), and medications are reviewed to determine perioperative continuation (e.g., beta-blockers, statins) or discontinuation (e.g., ACE inhibitors, anticoagulants, oral hypoglycemics, GLP-1 receptor agonists).

\medterm{Propofol}-- An alkylphenol derivative (2,6-diisopropylphenol) formulated as a 1\% or 2\% white, isotonic oil-in-water emulsion containing soybean oil, glycerol, and purified egg yolk phosphatide, serving as the most widely utilized intravenous hypnotic agent for the induction and maintenance of general anesthesia and procedural sedation. It acts primarily by allosterically binding to and activating $\beta$-subunits of the inhibitory GABA-A receptor complex, prolonging inhibitory postsynaptic chloride currents and suppressing neuronal excitability. Standard induction dose is 1.5 to 2.5 mg/kg IV in healthy adults (reduced to 1.0--1.5 mg/kg in elderly or compromised patients), producing smooth, rapid loss of consciousness in 30 to 45 seconds with a brief clinical duration (5 to 10 minutes) due to rapid redistribution into muscle and fat. Maintenance infusion rate for Total Intravenous Anesthesia (TIVA) is typically 100 to 200 mcg/kg/min (6 to 12 mg/kg/hr) or effect-site target-controlled infusion ($C_e$ 3 to 6 mcg/mL). It produces dose-dependent arterial and venous vasodilation and mild myocardial depression, characteristically causing a 20 to 30 percent reduction in mean arterial pressure, often blunting compensatory tachycardia by impairing the baroreceptor reflex. It causes profound, dose-dependent respiratory depression and blunts upper airway reflexes (facilitating laryngeal mask airway insertion without paralytics). It possesses potent antiemetic, antipruritic, and bronchodilating properties, and produces a rapid, clear-headed emergence. Pain on injection is common (mitigated by lidocaine pre-treatment or large antecubital veins). Propofol Infusion Syndrome (PRIS) is a rare, lethal complication associated with high-dose ($>4--5$ mg/kg/hr), prolonged ($>48$ hours) infusions, characterized by refractory metabolic acidosis, rhabdomyolysis, hyperkalemia, hepatomegaly, renal failure, and fatal bradyarrhythmic cardiac collapse.

\medterm{Pulse Oximetry}-- A non-invasive, continuous optical monitoring modality that measures arterial hemoglobin oxygen saturation ($SpO_2$) and pulse rate at a peripheral vascular bed (finger, toe, earlobe, forehead). Its operational principle combines dual-wavelength spectrophotometry with photoplethysmography based on the Beer-Lambert law: an optical sensor emits light at two distinct wavelengths—red light at 660 nm (which is absorbed predominantly by deoxygenated hemoglobin) and infrared light at 940 nm (absorbed predominantly by oxyhemoglobin). By analyzing the ratio of pulsatile (AC, arterial) light absorption to non-pulsatile (DC, tissue, venous blood, bone) background absorption, the microchip calculates functional arterial oxygen saturation. The standard pulse oximeter is highly accurate ($\pm 2\%$ between 70 and 100 percent saturation). Limitations and confounding factors: (1) Carboxyhemoglobin absorbs light at 660 nm identically to oxyhemoglobin, producing falsely elevated $SpO_2$ in carbon monoxide poisoning; (2) Methemoglobin absorbs equally at 660 and 940 nm, causing the calculated ratio to default toward 1.0, fixing the display reading at ~85\% regardless of actual arterial PaO2; (3) Severe peripheral vasoconstriction, hypothermia, shock, or severe hypovolemia diminishes the pulsatile signal, causing reading failure; (4) Intravascular dyes (methylene blue, indocyanine green) transiently cause severe falsely low $SpO_2$ readings; and (5) Severe anemia and motion artifact. Pulse oximetry measures oxygenation, not ventilation; normal SpO2 can persist during severe hypoventilation in patients receiving supplemental oxygen until sudden desaturation occurs.

\medterm{Rapid Sequence Induction}-- An intubation technique designed to minimize aspiration risk in patients with full stomach. Steps include preoxygenation with 100 percent oxygen for 3 to 5 minutes (8 vital capacity breaths or 3 minutes), application of cricoid pressure, administration of induction agent (propofol 2 mg/kg or etomidate 0.3 mg/kg) followed immediately by succinylcholine (1.5 mg/kg) or rocuronium (1.2 mg/kg), and intubation without bag-mask ventilation. If rocuronium is used, sugammadex (16 mg/kg) is available as a rescue reversal agent, providing rapid and complete reversal of rocuronium-induced neuromuscular blockade. RSI is indicated in full stomach, emergency surgery, obesity, pregnancy, gastrointestinal obstruction, and severe gastro-oesophageal reflux disease.

\medterm{Regional Anesthesia}-- The delivery of local anesthetic solutions near specific nerves, nerve plexuses, or the spinal cord to produce temporary, reversible loss of sensation and motor function over a discrete anatomical region without loss of consciousness. Subdivided into: (1) Neuraxial anesthesia (spinal and epidural blockade), where local anesthetics are placed within the subarachnoid space or epidural space to block spinal nerve roots; (2) Peripheral nerve and plexus blocks (e.g., brachial plexus, femoral, sciatic, and popliteal blocks), performed standardly with real-time ultrasound guidance and nerve stimulation to maximize block success and prevent inadvertent intraneural or intravascular injection; and (3) Truncal/fascial plane blocks (e.g., transversus abdominis plane [TAP], erector spinae plane [ESP], and rectus sheath blocks). Regional anesthesia provides superior postoperative analgesia, reduces systemic opioid consumption, minimizes nausea and sedation, facilitates early recovery, and attenuates perioperative immunosuppression. Contraindications include patient refusal, localized infection at the puncture site, uncorrected coagulopathy, severe hypovolemic shock, and pre-existing severe peripheral neuropathy in the distribution of the block.

\medterm{Remifentanil}-- A potent, ultra-short-acting synthetic opioid analgesic, chemically related to fentanyl, characterized by a unique ester linkage that renders it susceptible to rapid hydrolysis by non-specific plasma and tissue esterases. It has an extraordinarily short, context-sensitive half-time of approximately 3 to 5 minutes that remains independent of the duration of infusion, allowing immediate recovery and rapid cognitive emergence even after multi-hour infusions. Remifentanil is approximately 250 times more potent than morphine. Administered by continuous IV infusion (typically 0.05 to 0.5 mcg/kg/min) or target-controlled infusion (TCI, Minto model), it is widely utilized as the analgesic component of total intravenous anesthesia (TIVA) with propofol and for rigid neurosurgical, ENT, and cardiovascular procedures requiring precise hemodynamic control and immediate postoperative neurological assessment. Side effects include dose-dependent respiratory depression, Remifentanil-induced bradycardia and hypotension (especially with bolus administration $>1$ mcg/kg), chest wall rigidity, and rapid development of acute opioid-induced hyperalgesia (OIH) following discontinuation, which necessitates preemptive administration of longer-acting analgesics (e.g., morphine, hydromorphone, acetaminophen, NSAIDs, ketamine) prior to emergence.

\medterm{Rocuronium}-- An intermediate-acting, steroidal non-depolarizing neuromuscular blocking agent (NMBA) that serves as the most widely used paralytic in modern clinical anesthesia. It acts as a competitive antagonist at postjunctional nicotinic acetylcholine receptors at the neuromuscular junction, preventing acetylcholine binding and muscle depolarization without causing fasciculations. Dosing: Standard intubating dose for elective general anesthesia is 0.6 mg/kg IV (onset 60 to 90 seconds, clinical duration 30 to 45 minutes); for rapid sequence induction (RSI), a dose of 1.0 to 1.2 mg/kg IV produces ideal intubating conditions in 45 to 60 seconds (comparable to succinylcholine). Maintenance bolus dose is 0.15 mg/kg. It lacks significant active metabolites and is cleared primarily by hepatic uptake and biliary excretion (70\%) with secondary renal elimination (30\%). It possesses an excellent cardiovascular safety profile with no significant autonomic ganglion blockade, no muscarinic blockade, and no histamine release. A major clinical advantage is that its neuromuscular blockade can be immediately and completely reversed at any depth within 2 to 3 minutes by the selective encapsulating agent sugammadex.

\medterm{Ropivacaine}-- A pure S-(-)-enantiomer amino-amide local anesthetic developed to provide long-acting conduction blockade with a significantly wider safety margin than racemic bupivacaine. It reversibly blocks voltage-gated sodium channels in peripheral and neuraxial nerve fibers, preventing action potential propagation. Because of its slightly lower lipid solubility compared to bupivacaine, ropivacaine displays a greater degree of differentiation between sensory and motor nerve fibers (sensory-motor dissociation), producing dense sensory analgesia with significantly less pronounced and shorter-duration motor blockade at lower concentrations (0.1--0.2\%), making it the premier agent for labor epidural analgesia and postoperative peripheral nerve infusions where patient ambulation is prioritized. At higher concentrations (0.5--0.75\%), it provides dense surgical sensory and motor block for peripheral nerve and epidural anesthesia. Its pure S-(-)-stereochemical structure confers substantially reduced cardiotoxicity and central nervous system toxicity: it dissociates much more rapidly from cardiac sodium channels than R(+)-bupivacaine, making fatal ventricular arrhythmias much less likely. Maximum safe single dose is 3.0 mg/kg (up to 250 mg). Extensively metabolized in the liver via cytochrome P450 CYP1A2.

\medterm{Scavenging System}-- The anesthesia gas scavenging system collects and removes waste anesthetic gases from the operating room environment to protect staff from chronic exposure. Active scavenging uses negative pressure from vacuum source; passive scavenging vents to exhaust through gravity. The interface between the anesthesia machine and scavenging system includes a positive pressure relief valve (to prevent barotrauma) and negative pressure relief valve (to prevent excessive suction). Occupational exposure limits for nitrous oxide are 25 ppm (time-weighted average over 8 hours) and for volatile agents 0.5-2 ppm depending on the agent. Proper scavenging is a critical occupational health and safety requirement in all anaesthetising locations, with periodic environmental monitoring to verify adequate system function and compliance with exposure limits.

\medterm{Sedation Levels}-- A clinically defined pharmacological continuum of central nervous system depression established by the American Society of Anesthesiologists (ASA), ranging from minimal sedation to general anesthesia: (1) Minimal Sedation (Anxiolysis): A drug-induced state during which patients respond normally to verbal commands; cognitive function and physical coordination may be impaired, but airway reflexes, ventilatory function, and cardiovascular stability are unaffected. (2) Moderate Sedation/Analgesia ('Conscious Sedation'): A drug-induced depression of consciousness during which patients respond purposefully to verbal commands, either alone or accompanied by light tactile stimulation; no interventions are required to maintain a patent airway, and spontaneous ventilation and cardiovascular function are adequate. (3) Deep Sedation/Analgesia: A drug-induced depression of consciousness during which patients cannot be easily aroused but respond purposefully following repeated or painful stimulation; patients may require assistance in maintaining a patent airway, spontaneous ventilation may be inadequate, but cardiovascular function is usually maintained. (4) General Anesthesia: A drug-induced loss of consciousness during which patients are unarousable, even by painful stimulation; the ability to independently maintain ventilatory function is often impaired, patients often require assistance in maintaining a patent airway, positive-pressure ventilation may be required, and cardiovascular function may be impaired. Because sedation is a dynamic continuum, practitioners must be qualified and equipped to rescue patients who drift into a level of sedation deeper than initially intended.

\medterm{Sevoflurane}-- A fluorinated methyl isopropyl ether that is the most widely utilized volatile inhalational anesthetic agent in modern clinical practice. It has an pleasant, non-pungent, sweet odor and causes minimal airway irritation, making it the universally preferred agent for inhalational mask induction in pediatric patients and adult difficult airways. Its minimum alveolar concentration (MAC) is approximately 2.0\% in pure oxygen (and 1.1\% in 65\% nitrous oxide) in young adults, decreasing with age, hypothermia, and opioid administration. It possesses a low blood-gas partition coefficient of 0.65, which allows rapid alveolar wash-in, prompt induction of anesthesia, and swift emergence following discontinuation. It undergoes approximately 3 to 5 percent hepatic metabolism by CYP2E1, releasing inorganic fluoride ions, but clinical fluoride-induced nephrotoxicity has not been observed. When exposed to carbon dioxide absorbents containing strong alkaline bases (barium hydroxide or potassium/sodium hydroxide), sevoflurane undergoes degradation to fluoromethyl-2,2-difluoro-1-(trifluoromethyl)vinyl ether (Compound A), which is nephrotoxic in laboratory rats; clinical guidelines mandate maintaining fresh gas flow at $\ge 1$ to 2 L/min during sevoflurane administration to prevent Compound A accumulation. Like all volatile agents, sevoflurane produces dose-dependent reductions in systemic vascular resistance and blood pressure, suppresses ventilation, increases cerebral blood flow at concentrations $>1$ MAC, and is an absolute trigger for malignant hyperthermia.

\medterm{Spinal Analgesia}-- Intrathecal administration of a local anesthetic and/or opioid primarily to provide analgesia, often with limited motor block. Intrathecal opioids can extend postoperative pain relief but may cause pruritus, urinary retention, respiratory depression, and other opioid effects. It should be distinguished from spinal anesthesia, in which the intrathecal local anesthetic dose is intended to produce surgical sensory and motor blockade.

\medterm{Spinal Anesthesia}-- A neuraxial regional anesthetic technique (subarachnoid block) accomplished by injecting a small volume of local anesthetic solution directly into the cerebrospinal fluid within the lumbar subarachnoid space. The lumbar puncture is performed strictly below the termination of the adult spinal cord (conus medullaris, terminating at L1--L2) at the L3--L4 or L4--L5 interspace, with the patient in the sitting or lateral decubitus position. The needle (preferably a 25--27 gauge pencil-point needle like Whitacre or Sprotte to minimize post-dural puncture headache) traverses skin, subcutaneous fat, supraspinous ligament, interspinous ligament, ligamentum flavum, epidural space, and dura-arachnoid membrane, confirmed by the free flow of clear CSF. Hyperbaric solutions (local anesthetic prepared with dextrose to increase density relative to CSF, such as 0.75\% bupivacaine in 8.25\% dextrose) or isobaric solutions are administered. Produces immediate, dense sensory, motor, and sympathetic blockade. Target dermatomal levels: T4 (nipple line) for cesarean section, T10 (umbilicus) for lower abdominal, hip, and TURP surgery, and L1 for lower extremity surgery. Physiological effects: Rapid sympathetic preganglionic blockade causes extensive venodilation and arterial vasodilation, resulting in marked hypotension; if the block ascends to T1--T4 (blocking cardioaccelerator fibers), profound bradycardia ensues (treated with fluids, ephedrine, phenylephrine, or epinephrine). Complications include high/total spinal anesthesia, post-dural puncture headache (PDPH), transient neurological symptoms (TNS), spinal hematoma, and meningitis.

\medterm{Succinylcholine}-- A quaternary ammonium compound consisting of two joined acetylcholine molecules (diacetylcholine) that acts as the only clinically utilized depolarizing neuromuscular blocking agent. It binds to nicotinic acetylcholine receptors at the motor endplate, opening receptor channels and producing sustained motor endplate depolarization. This initial stimulation produces characteristic transient, disorganized muscle contractions known as fasciculations (Phase I block, showing lack of fade on train-of-four stimulation), followed immediately by flaccid neuromuscular paralysis because voltage-gated sodium channels in the perijunctional membrane remain inactivated. Standard intubation dose is 1.0 to 1.5 mg/kg IV (onset 30 to 60 seconds, clinical duration 5 to 10 minutes), making it the historical gold standard for rapid sequence induction (RSI) in patients with a full stomach or high aspiration risk. It is rapidly hydrolyzed in the bloodstream by plasma pseudocholinesterase (butyrylcholinesterase); in patients with genetic pseudocholinesterase deficiency (homozygous atypical enzyme, identified by a low dibucaine number), apnea may be prolonged for several hours. Adverse effects and contraindications: (1) Acute hyperkalemic cardiac arrest: causes an obligate serum potassium release of ~0.5 mEq/L, which can be massive and fatal in patients with denervation injuries (spinal cord injury, stroke), major burns, muscular dystrophies, massive crush injury, and severe sepsis due to upregulation of extrajunctional acetylcholine receptors; (2) Malignant hyperthermia trigger; (3) Severe sinus bradycardia and nodal rhythms (especially with repeated doses or in pediatric patients, prevented with atropine); (4) Masseter muscle spasm; (5) Postoperative myalgias; and (6) Transient increases in intraocular, intracranial, and intragastric pressure.

\medterm{Sugammadex}-- A modified gamma-cyclodextrin derivative that acts as a selective relaxant binding agent, designed specifically to reverse the neuromuscular blockade induced by steroidal neuromuscular blocking agents (predominantly rocuronium, and to a lesser extent vecuronium). Its lipophilic core encapsulates the steroidal guest molecule in a 1:1 intermolecular complex with high affinity, rapidly decreasing free plasma concentration and drawing the relaxant away from the nicotinic acetylcholine receptors at the neuromuscular junction into the vascular compartment for renal excretion. Dosing is determined strictly by the monitored depth of neuromuscular blockade: (1) 2 mg/kg IV for the reversal of moderate blockade, defined as the reappearance of at least the second twitch ($T_2$) of the train-of-four (TOF) response; (2) 4 mg/kg IV for deep blockade, defined as 1 to 2 post-tetanic counts (PTC) with 0 twitches on TOF; and (3) 16 mg/kg IV for immediate rescue reversal approximately 3 minutes following a single intubating dose of 1.2 mg/kg rocuronium (e.g., in a cannot intubate, cannot oxygenate emergency). Unlike anticholinesterase agents (neostigmine), sugammadex acts independently of acetylcholinesterase, producing rapid, complete reversal without cholinergic side effects (eliminating the need for co-administered anticholinergic agents like glycopyrrolate) and without risk of recurarization when dosed properly. Adverse effects include hypersensitivity/anaphylaxis (approximately 1 in 2,500), transient bradycardia (which can rarely progress to cardiac arrest, responding to atropine), and binding of endogenous and exogenous steroids, notably oral contraceptives (patients must be advised to use alternative non-hormonal contraception for 7 days following administration).

\medterm{Supraglottic Airway Devices}-- Extraglottic airway management devices (SADs) that sit outside the trachea, positioning an inflatable or anatomically contoured cuff directly above the glottis to provide an airtight seal over the laryngeal inlet. First-generation SADs (e.g., classic Laryngeal Mask Airway [LMA Classic], Flexible LMA) provide a basic airway seal (seal pressures 15--20 cmH2O) suitable for spontaneously breathing or low-pressure ventilated patients, but lack protection against gastric regurgitation. Second-generation SADs (e.g., LMA ProSeal, LMA Supreme, i-gel, Ambu AuraGain) represent modern clinical standards: they incorporate a dedicated gastric drainage tube (allowing decompression of the stomach and venting of regurgitated gastric contents), an integral bite block, and higher mucosal seal pressures ($>25--30$ cmH2O), enabling safe positive-pressure ventilation. SADs are fundamental to difficult airway rescue algorithms as the primary conduit for oxygenation when face-mask ventilation or direct endotracheal intubation fails, and can serve as guides for fiberoptic-assisted intubation. Contraindications include severe gastroesophageal reflux, morbid obesity with high airway pressures, delayed gastric emptying/full stomach, pharyngeal pathology, and situations requiring high ventilatory pressures.

\medterm{Target-Controlled Infusion}-- A computerized intravenous infusion delivery system that utilizes population-based three-compartment pharmacokinetic and pharmacodynamic mathematical models to achieve and maintain a user-selected target drug concentration in plasma ($C_p$) or effect-site ($C_e$). The infusion pump microprocessor automatically calculates and continuously adjusts the variable infusion rate based on patient covariates (age, weight, height, sex) to rapidly attain the desired concentration with an initial bolus and then balance drug elimination and redistribution. Common clinical models include the Marsh model (plasma-targeted) and Schnider model (effect-site targeted) for propofol, and the Minto model for remifentanil. TCI provides superior hemodynamic stability, reduces anesthetic consumption, prevents unintended awareness and drug accumulation, and allows predictable emergence times compared to manual infusion techniques. Limitations include interindividual pharmacokinetic variability in critically ill, elderly, or pediatric patients, and the requirement for dedicated TCI-compatible infusion hardware.

\medterm{Total Intravenous Anesthesia}-- Total intravenous anesthesia (TIVA) uses intravenous agents exclusively for induction and maintenance of anesthesia, avoiding volatile inhalation agents. The most common TIVA regimen combines propofol infusion (50 to 200 mcg/kg/min) with remifentanil infusion (0.1 to 0.5 mcg/kg/min). TIVA advantages include rapid emergence, reduced postoperative nausea and vomiting, suitability for patients with malignant hyperthermia risk, and use in procedures requiring volatile anesthetic avoidance. TIVA disadvantages include the need for dedicated infusion pumps, risk of propofol infusion syndrome with prolonged high-dose infusions, and the lack of an end-tidal volatile anesthetic concentration monitor to confirm anesthetic depth; BIS or entropy monitoring is typically used with TIVA to assess depth of anesthesia and reduce the risk of intraoperative awareness.

\medterm{Train-of-Four Monitoring}-- Train-of-four monitoring assesses neuromuscular blockade by delivering four supramaximal electrical stimuli at 2 Hz to a peripheral nerve, commonly the ulnar nerve at the wrist. The ratio of the fourth twitch height to the first twitch height estimates recovery from non-depolarizing blockade. A quantitative TOF ratio of at least 0.9 is generally required before extubation to reduce residual neuromuscular blockade and associated respiratory complications.

\medterm{Tranexamic Acid}-- An antifibrinolytic agent that inhibits lysine-binding sites on plasminogen and reduces clot breakdown. It is used in selected trauma, obstetric, cardiac, orthopedic, and bleeding indications according to current protocols. Dose and route depend on the indication and renal function. Important adverse effects include seizures, hypersensitivity, and potential thrombotic complications; contraindications and precautions must be assessed clinically rather than inferred from a history of every prior thromboembolic event.

\medterm{Vasopressors in Anesthesia}-- Vasopressors maintain blood pressure during anesthesia-induced vasodilation. Phenylephrine is a pure alpha-1 agonist causing vasoconstriction without cardiac effects, used for pure vasodilatory hypotension. Dose: bolus 50 to 100 mcg IV or infusion 0.05 to 0.5 mcg/kg/min. Ephedrine stimulates alpha and beta receptors, increasing both heart rate and blood pressure. Dose: 5 to 10 mg IV bolus. Norepinephrine is a potent alpha-1 and beta-1 agonist, used for refractory hypotension and septic shock. Dose: 0.05 to 0.5 mcg/kg/min infusion. Dopamine acts on dopamine, beta, and alpha receptors in a dose-dependent manner: low dose (1-5 mcg/kg/min) mainly acts on dopamine receptors dilating renal and mesenteric vessels; intermediate dose (5-10 mcg/kg/min) acts on beta-1 receptors increasing contractility and heart rate; high dose (>10 mcg/kg/min) acts on alpha-1 receptors causing vasoconstriction.

\medterm{Video Laryngoscopy}-- Video laryngoscopy uses a camera-equipped blade to display an indirect view of the glottis on a screen. It can improve glottic visualization in many patients, including some with limited neck movement or difficult airway anatomy, but it does not replace a difficult-airway plan, effective oxygenation, or preparation for alternative techniques. The operator must understand the device-specific blade, tube-delivery, and rescue techniques.

\medterm{Volatile Anesthetics}-- A group of halogenated hydrocarbon and ether liquids (primarily sevoflurane, desflurane, and isoflurane; historically halothane and enflurane) that are vaporized into gas form and inhaled to induce and maintain general anesthesia. Their primary molecular mechanism involves the allosteric potentiation of inhibitory ligand-gated ion channels (GABA-A and glycine receptors) in the brain and spinal cord, alongside the inhibition of excitatory receptors (nicotinic acetylcholine, NMDA, AMPA, and two-pore domain potassium channels), producing hypnosis, amnesia, and skeletal muscle immobility in response to noxious stimulation. Potency is measured by Minimum Alveolar Concentration (MAC), which correlates inversely with lipid solubility (Meyer-Overton correlation). Pharmacokinetics are governed by tissue and blood solubilities: low blood-gas partition coefficients (desflurane 0.42, sevoflurane 0.65) enable rapid induction and emergence, whereas higher solubility (isoflurane 1.4) slows equilibration. Physiological effects: Dose-dependent depression of myocardial contractility, reduction in systemic vascular resistance (causing hypotension), blunting of ventilatory drive and hypoxic pulmonary vasoconstriction (HPV), bronchodilation, and cerebral vasodilation (which increases cerebral blood flow and intracranial pressure at $>1$ MAC). All volatile anesthetics are potent triggers of malignant hyperthermia in genetically susceptible individuals and may contribute to postoperative nausea and vomiting.

